% Generated by scripts/build_reports.py; edit reviewed Markdown inputs.
% Compile with XeLaTeX.
% Options for packages loaded elsewhere
\PassOptionsToPackage{unicode}{hyperref}
\PassOptionsToPackage{hyphens}{url}
\PassOptionsToPackage{dvipsnames,svgnames,x11names}{xcolor}
\documentclass[
  11pt,
]{article}
\usepackage{xcolor}
\usepackage[margin=1in]{geometry}
\usepackage{amsmath,amssymb}
\setcounter{secnumdepth}{-\maxdimen} % remove section numbering
\usepackage{iftex}
\ifPDFTeX
  \usepackage[T1]{fontenc}
  \usepackage[utf8]{inputenc}
  \usepackage{textcomp} % provide euro and other symbols
\else % if luatex or xetex
  \usepackage{unicode-math} % this also loads fontspec
  \defaultfontfeatures{Scale=MatchLowercase}
  \defaultfontfeatures[\rmfamily]{Ligatures=TeX,Scale=1}
\fi
\usepackage{lmodern}
\ifPDFTeX\else
  % xetex/luatex font selection
\fi
% Use upquote if available, for straight quotes in verbatim environments
\IfFileExists{upquote.sty}{\usepackage{upquote}}{}
\IfFileExists{microtype.sty}{% use microtype if available
  \usepackage[]{microtype}
  \UseMicrotypeSet[protrusion]{basicmath} % disable protrusion for tt fonts
}{}
\makeatletter
\@ifundefined{KOMAClassName}{% if non-KOMA class
  \IfFileExists{parskip.sty}{%
    \usepackage{parskip}
  }{% else
    \setlength{\parindent}{0pt}
    \setlength{\parskip}{6pt plus 2pt minus 1pt}}
}{% if KOMA class
  \KOMAoptions{parskip=half}}
\makeatother
\usepackage{longtable,booktabs,array}
\usepackage{caption}
\captionsetup[table]{skip=6pt}
\newcounter{none} % for unnumbered tables
\usepackage{calc} % for calculating minipage widths
% Correct order of tables after \paragraph or \subparagraph
\usepackage{etoolbox}
\makeatletter
\patchcmd\longtable{\par}{\if@noskipsec\mbox{}\fi\par}{}{}
\makeatother
% Allow footnotes in longtable head/foot
\IfFileExists{footnotehyper.sty}{\usepackage{footnotehyper}}{\usepackage{footnote}}
\makesavenoteenv{longtable}
\setlength{\emergencystretch}{3em} % prevent overfull lines
\providecommand{\tightlist}{%
  \setlength{\itemsep}{0pt}\setlength{\parskip}{0pt}}
\usepackage{amsmath,amssymb}
\usepackage{xurl}
\usepackage{fancyhdr}
\pagestyle{fancy}
\fancyhf{}
\fancyhead[L]{\small Pure Thought | Revised research specifications}
\fancyfoot[L]{\small 2026-09-17}
\fancyfoot[R]{\small\thepage}
\setlength{\headheight}{15pt}
\setlength{\emergencystretch}{3em}
\allowdisplaybreaks
\usepackage{bookmark}
\IfFileExists{xurl.sty}{\usepackage{xurl}}{} % add URL line breaks if available
\urlstyle{same}
% fallback for those not using the hyperref driver hyperxmp:
\makeatletter
\@ifundefined{xmpquote}{\newcommand{\xmpquote}[1]{#1}}{}
\makeatother
\hypersetup{
  pdftitle={Thirty synthesis challenges: revised summaries},
  pdfauthor={Pure Thought research portfolio},
  colorlinks=true,
  linkcolor={blue},
  filecolor={Maroon},
  citecolor={Blue},
  urlcolor={blue},
  pdfcreator={LaTeX via pandoc}}

\title{Thirty synthesis challenges: revised summaries}
\author{Pure Thought research portfolio}
\date{Revised 17 September 2026}

\begin{document}
\maketitle
\thispagestyle{empty}
\clearpage

{
\setcounter{tocdepth}{1}
\tableofcontents
\clearpage
}
\section{Thirty synthesis challenges: revised problem summaries}

Updated 2026-09-17. This synthesis retains the original synthesis
numbering. It is a different portfolio from the canonical 30 PRDs. Every
topic is mapped to its revised destination; the 12 distinct supplemental
objectives are preserved. Claims and milestones follow the updated
review and evidence standard.

\href{https://github.com/rick-stevens-ai/Pure-Thought/blob/main/COMPLETE-PROBLEM-SUMMARIES.md}{All
42 complete summaries} \ensuremath{\cdot}
\href{https://github.com/rick-stevens-ai/Pure-Thought/blob/main/reviews/SYNTHESIS-CROSSWALK.md}{Full
crosswalk} \ensuremath{\cdot}
\href{https://github.com/rick-stevens-ai/Pure-Thought/blob/main/NOTES/30-Pure-Thought-Challenges-Report.pdf}{Revised
synthesis PDF} \ensuremath{\cdot}
\href{https://github.com/rick-stevens-ai/Pure-Thought/blob/main/archive/original-synthesis/30-Most-Compelling-Pure-Thought-AI-Challenges.md}{Preserved
original}

\section{\texorpdfstring{Synthesis 01. AdS\ensuremath{{}_{3}} ``Pure
Gravity'' via the Modular
Bootstrap}{Synthesis 01.  ``Pure Gravity'' via the Modular Bootstrap}}

\textbf{Mapping:}
\href{https://github.com/rick-stevens-ai/Pure-Thought/blob/main/PRDs/01-AdS3-Modular-Bootstrap.md}{PRD
01}. Covered by canonical revision.

\subsection{01. Extremal holomorphic CFT partition functions: certified
necessary tests}

\textbf{Question.} For c=24k, which extremal holomorphic
partition-function candidates pass exact genus-one spectral tests, and
can any candidate be excluded by an additional, explicitly stated
consistency condition?

\textbf{Why it matters.} Quantum gravity is hard partly because we do
not know which mathematically plausible theories are internally
consistent. These tests can eliminate candidates without experiments or
a complete construction. They also connect gravity to the arithmetic of
modular functions: counting quantum states becomes a problem about exact
number sequences.

\textbf{Scope and model.} Start with bosonic holomorphic theories, a
unique vacuum, c=24k, and k=1,\ldots,4. This is a deliberately
restricted version of the original gravity question. Use h for
holomorphic weight; do not identify it with the full scaling dimension
\ensuremath{\Delta}=h+\ensuremath{\bar{h}}. The extremal ansatz has no
nonvacuum Virasoro primary with h\ensuremath{\le}k. Treat the nonchiral
modular bootstrap as a separate extension with independent
(h,\ensuremath{\bar{h}}) data.

Keep the motivating existence question, but make the first deliverable a
necessary-condition test. A modular function with integer coefficients
is not a constructed CFT or a demonstrated gravity dual.

\subsubsection{Mathematical formulation and claim boundary}

\href{https://github.com/rick-stevens-ai/Pure-Thought/blob/main/PRDs/01-AdS3-Modular-Bootstrap.md}{Canonical
specification}.

For the holomorphic sector with \(c=24k\), \(q=e^{2\pi i\tau}\) and
\(J=j-744\),

\[\chi_0=q^{-k}\prod_{n=2}^{\infty}(1-q^n)^{-1},\qquad
\chi_h=\frac{q^{h-k+1/24}}{\eta(\tau)}\quad(h>0).\]

\[Z_k=P_k(J)=\chi_0+\sum_{h\ge k+1}d_h\chi_h,\qquad
P_k(J)-\chi_0=O(q).\]

Determine \(P_k\) exactly and test \(d_h\in\mathbb Z_{\ge0}\) through a
declared \(H\). At \(k=1\), \(d_2=196883\). Passing this finite
necessary test does not construct the CFT or establish OPE consistency.

\textbf{Deliverables.} Export the polynomial coefficients, exact
q-series, primary multiplicities, conventions, cutoff, and a checker
that independently expands \ensuremath{\eta} and J. An exclusion must
name the violated necessary condition. An additional genus-two or OPE
test must state its own assumptions and how its constraints follow from
CFT consistency.

\textbf{Validation.} Recover J=q\ensuremath{{}^{-1}}+196884q+… at k=1
and d\ensuremath{{}_{2}}=196883 after subtracting the vacuum descendant.
Test the checker with a deliberately altered coefficient. Check
agreement of two series-construction routes. Do not report all-order
coefficient positivity without a tail theorem.

\textbf{Milestones.} Baseline: exact k=1,\ldots,4 tables with explicitly
finite coverage. Strong: one additional sewing/OPE condition with
independently checked consequences. Research extension: a new exclusion
theorem or a full CFT construction; neither is a required outcome or a
promised deadline.

\href{https://github.com/rick-stevens-ai/Pure-Thought/blob/main/reviews/01-AdS3-Modular-Bootstrap.md}{Detailed
critique}.

\section{Synthesis 02. Gravitational Positivity \& Causality Bounds
(Flat Space)}

\textbf{Mapping:}
\href{https://github.com/rick-stevens-ai/Pure-Thought/blob/main/PRDs/02-Gravitational-Positivity-Bounds.md}{PRD
02}. Covered by canonical revision.

\subsection{02. Conditional positivity bounds for a specified
gravitational amplitude}

\textbf{Question.} Under a fully stated infrared prescription and
high-energy assumption, what certified inequality constrains one
selected low-energy coefficient of scalar scattering coupled to gravity?

\textbf{Why it matters.} These bounds can tell us that an apparently
sensible low-energy model cannot fit particular assumptions about its
underlying physics. That is valuable when direct access to microscopic
gravity is unavailable. A reliable bound also prevents researchers from
interpreting an artifact of a singular calculation as a physical
restriction.

\textbf{Scope and model.} Use four-dimensional scattering of one
identical real massive scalar coupled to Einstein gravity. This replaces
the redundant schematic pure-gravity R\ensuremath{{}^{2}} coefficient
list with an amplitude-defined observable. Set m=1 as the mass unit;
specify \ensuremath{\kappa}, the EFT subtraction scale, perturbative
order, and whether the target is a tree coefficient or a renormalized
coefficient. Graviton helicity amplitudes are a later extension.

The first task is to establish a valid sum rule. Do not run an
optimization until the pole subtraction, massless cuts, crossed channel,
and large-contour contribution have been derived.

\subsubsection{Mathematical formulation and claim boundary}

\href{https://github.com/rick-stevens-ai/Pure-Thought/blob/main/PRDs/02-Gravitational-Positivity-Bounds.md}{Canonical
specification}.

\[s+t+u=4m^2,\qquad \nu=s-2m^2+\frac t2,\qquad
b_2(t)=\frac12\partial_\nu^2 B(0,t).\]

\[B=A-A_{\mathrm{light\ poles}}-A_{\mathrm{chosen\ low\ cuts}},\qquad
b_2(t)=I_{\mathrm{UV}}(t)+R_\infty(t).\]

The last equality is the sum rule to derive, with its kernel,
subtractions and contour assumptions. Only a justified smearing or
regulated forward prescription can turn it into a positivity statement.
If that prescription proves \(I_{\mathrm{UV}}\ge0\) and
\(R_\infty\ge-R_*\), then \(b_2\ge-R_*\). An uncontrolled contour term
must remain unknown.

\textbf{Deliverables.} Return a derivation dossier, normalized
coefficient definition, assumptions, and either a conditional inequality
b\ensuremath{{}_{2}}\ensuremath{\ge}L or an explanation of the missing
hypothesis. If optimizing, export a dual functional with certified
kernel positivity and all truncation errors. A surviving coefficient is
``not excluded by these tests,'' not ``UV complete.''

\textbf{Validation.} Check s\ensuremath{\leftrightarrow}u symmetry,
dimensional consistency, and pole residues. Reproduce the
nongravitational forward positivity argument in the
\ensuremath{\kappa}\ensuremath{\to}0 limit. Use a fully specified
example amplitude to test the subtraction algebra. Perturb a certificate
beyond its margin and verify rejection.

\textbf{Milestones.} Baseline: one correct sum rule and its gravity-free
limit. Strong: a numerically useful certified conditional bound with a
finite positive margin. Research extension: improve the bound or relax a
high-energy assumption; label any sharpness claim with both matching
constructions and upper/lower evidence.

\href{https://github.com/rick-stevens-ai/Pure-Thought/blob/main/reviews/02-Gravitational-Positivity-Bounds.md}{Detailed
critique}.

\section{Synthesis 03. Celestial CFT Bootstrap}

\textbf{Mapping:}
\href{https://github.com/rick-stevens-ai/Pure-Thought/blob/main/PRDs/03-Celestial-CFT-Bootstrap.md}{PRD
03}. Covered by canonical revision.

\subsection{03. Celestial amplitudes: distributional crossing and
soft-limit consistency}

\textbf{Question.} Can a regulated celestial transform of a specified
four-graviton tree amplitude reproduce Lorentz covariance, crossing, and
the leading conformally soft residue in one consistent convention?

\textbf{Why it matters.} A successful reformulation could expose
symmetries or constraints that are difficult to see in ordinary
scattering variables. That would give theorists another way to
investigate gravity. But a change of mathematical language is useful
only if singularities, conservation laws and limiting cases survive the
translation correctly.

\textbf{Scope and model.} Fix four-dimensional Einstein gravity at tree
level, a nonzero four-graviton helicity amplitude, incoming/outgoing
signs \ensuremath{\epsilon}\ensuremath{{}_{i}}, and
p\ensuremath{{}_{i}}=\ensuremath{\epsilon}\ensuremath{{}_{i}}\ensuremath{\omega}\ensuremath{{}_{i}}q(z\ensuremath{{}_{i}},\ensuremath{\bar{z}}\ensuremath{{}_{i}}).
Declare metric, helicity, normalization, and analytic-continuation
conventions. The target is an identity pipeline for amplitudes, not a
classification of celestial CFTs.

Use
\ensuremath{\Delta}\ensuremath{{}_{i}}=1+i\ensuremath{\lambda}\ensuremath{{}_{i}}
on the principal series where appropriate and
h\ensuremath{{}_{i}}=(\ensuremath{\Delta}\ensuremath{{}_{i}}+J\ensuremath{{}_{i}})/2,
\ensuremath{\bar{h}}\ensuremath{{}_{i}}=(\ensuremath{\Delta}\ensuremath{{}_{i}}\ensuremath{-}J\ensuremath{{}_{i}})/2.
Identify explicitly which continuations away from that contour are used
for soft residues.

\subsubsection{Mathematical formulation and claim boundary}

\href{https://github.com/rick-stevens-ai/Pure-Thought/blob/main/PRDs/03-Celestial-CFT-Bootstrap.md}{Canonical
specification}.

\[p_i=\epsilon_i\omega_i q_i,\quad
\Delta_i=1+i\lambda_i,\quad
h_i=\frac{\Delta_i+J_i}{2},\quad
\bar h_i=\frac{\Delta_i-J_i}{2}.\]

\[\widetilde{\mathcal A}=
\left(\prod_{i=1}^4\int_0^\infty d\omega_i\,\omega_i^{\Delta_i-1}\right)
\delta^{(4)}\!\left(\sum_i\epsilon_i\omega_iq_i\right)\mathcal M_4.\]

The expression is a regulated distribution. A leading soft term
proportional to \(\omega_s^{-1}\) yields the local integral
\(\int_0^\Lambda\omega_s^{\Delta_s-2}d\omega_s=\Lambda^{\Delta_s-1}/(\Delta_s-1)\)
before continuation, exposing the pole at \(\Delta_s=1\). Support,
helicity, crossing and regulator removal remain proof obligations;
ordinary Euclidean OPE positivity is not assumed.

\textbf{Deliverables.} Produce exact momentum-space expressions,
regulated transform formulas, support conditions, and symbolic
covariance/crossing/soft identities. If a transform cannot be defined
under the chosen assumptions, return a precise obstruction at that step.
A positive bootstrap is an extension only after deriving an applicable
inner product and positivity statement.

\textbf{Validation.} Check the original amplitude's factorization first.
Test Mellin formulas against analytic one-variable examples in their
convergence strips. Verify soft residues before and after transformation
under stated limit-interchange hypotheses. Use smeared distributions for
numerical checks; never evaluate a delta distribution as an ordinary
function.

\textbf{Milestones.} Baseline: one four-point transform with explicit
support and checked leading soft residue. Strong: a second
helicity/process example and a verified crossing relation. Research
extension: derive a justified positive constraint in a restricted
subsector before attempting an SDP island.

\href{https://github.com/rick-stevens-ai/Pure-Thought/blob/main/reviews/03-Celestial-CFT-Bootstrap.md}{Detailed
critique}.

\section{Synthesis 04. Modular-Lightcone Bootstrap for Holographic CFTs}

\textbf{Mapping:}
\href{https://github.com/rick-stevens-ai/Pure-Thought/blob/main/PRDs/04-Modular-Lightcone-Bootstrap.md}{PRD
04}. Covered by canonical revision.

\subsection{04. Lightcone bootstrap with controlled large-spin
remainders}

\textbf{Question.} For one scalar four-point function in a unitary CFT
with d\textgreater2, what controlled bounds on large-spin double-twist
data follow from specified low-twist exchanges?

\textbf{Why it matters.} Consistency can reveal features shared by many
strongly interacting theories, even when their detailed dynamics cannot
be solved. Large-spin patterns are also useful in theories proposed to
describe gravity through a lower-dimensional boundary. Controlled errors
turn suggestive asymptotic patterns into statements that can actually
constrain a finite calculation.

\textbf{Scope and model.} Start in d=3 with an identical scalar
\ensuremath{\phi} of fixed dimension
\ensuremath{\Delta}\ensuremath{\phi}. Specify the normalization of C\_T,
the stress-tensor OPE coefficient fixed by its Ward identity, and an
explicit twist gap above the low-twist exchanges retained. Use a
large-C\_T expansion only when a factorization assumption and its error
order are stated.

This challenge studies analytic large-spin behavior. Challenge 07
instead studies the spinning stress-tensor correlator. Torus modular
invariance is restricted to a separate d=2 variant.

\subsubsection{Mathematical formulation and claim boundary}

\href{https://github.com/rick-stevens-ai/Pure-Thought/blob/main/PRDs/04-Modular-Lightcone-Bootstrap.md}{Canonical
specification}.

\[F_{\Delta,\ell}=v^{\Delta_\phi}g_{\Delta,\ell}(u,v)
-u^{\Delta_\phi}g_{\Delta,\ell}(v,u),\qquad
F_{\mathrm{id}}+\sum_{\mathcal O}\lambda_{\phi\phi\mathcal O}^{2}F_{\Delta,\ell}=0.\]

\[\Delta_{n,\ell}=2\Delta_\phi+2n+\ell+\gamma_{n,\ell},\qquad
|\gamma_{0,\ell}-\widehat\gamma_{0,\ell}|\le R(\ell)\quad(\ell\ge\ell_0).\]

The inequality is a target requiring explicit spectral and Regge
assumptions and a proved remainder. An asymptotic series alone does not
imply this finite-spin inequality. Higher-dimensional crossing is not
torus modular invariance.

\textbf{Deliverables.} Export block normalizations, low-twist input, the
derivation or inversion formula, coefficients, and a decomposition of R
into omitted exchanges, expansion terms, and numerical errors. State
clearly whether the result is an identity, a conditional bound, or an
asymptotic estimate.

\textbf{Validation.} Use generalized-free-field crossing as an algebraic
benchmark, while noting that it is not by itself the target local CFT
with a finite stress tensor. Recover double-twist accumulation in a
known analytic example. Verify every uniform estimate on the specified
cross-ratio domain and all spins covered by the claim.

\textbf{Milestones.} Baseline: reproduce an established leading
large-spin result with conventions checked. Strong: obtain one explicit
remainder estimate, or identify the extra data needed to obtain it.
Research extension: a tighter conditional bound on low-twist exchanges
or bulk couplings with a stated holographic dictionary.

\href{https://github.com/rick-stevens-ai/Pure-Thought/blob/main/reviews/04-Modular-Lightcone-Bootstrap.md}{Detailed
critique}.

\section{Synthesis 05. Positive Geometry for Gravity}

\textbf{Mapping:}
\href{https://github.com/rick-stevens-ai/Pure-Thought/blob/main/PRDs/05-Positive-Geometry-Gravity.md}{PRD
05}. Covered by canonical revision.

\subsection{05. Positive-geometry tests for a fixed gravity integrand}

\textbf{Question.} For a specified low-loop supergravity integrand, does
a precisely bounded class of candidate positive geometries reproduce its
differential form, including every boundary and pole at infinity?

\textbf{Why it matters.} A geometric organization could explain
cancellations and suggest more efficient calculations or new identities
in quantum gravity. Even a carefully limited impossibility theorem is
informative: it tells us which geometric proposals are too restrictive
and where a different idea is needed.

\textbf{Scope and model.} Start with the four-point one-loop N=8
supergravity integrand, in a declared four-dimensional representation
and helicity sector. Record any regulator and the limits of
four-dimensional cut reconstruction. Derive or reproduce the target
before searching. Tree-level and scalar-box canonical forms are
calibration examples.

Define the search class before claiming an obstruction: ambient
coordinates, dimension, allowed rational maps, inequalities, permitted
external prefactors, and bounds on degree/number of boundaries. ``No
geometry exists'' is not a finite computational question without such
restrictions.

\subsubsection{Mathematical formulation and claim boundary}

\href{https://github.com/rick-stevens-ai/Pure-Thought/blob/main/PRDs/05-Positive-Geometry-Gravity.md}{Canonical
specification}.

\[\Omega_{\mathrm{target}}=I(\ell;k)\,d^4\ell,
\qquad \Omega_{\mathrm{target}}\stackrel{?}{=}f(k)\,\pi_*\Omega_G.\]

For a standard positive geometry, each oriented boundary must satisfy
the canonical residue relation
\(\operatorname{Res}_{\partial G}\Omega_G=\Omega_{\partial G}\). Fix the
geometry class, maps, prefactors, representation and equality convention
first. Verify finite boundaries, intersections and infinity; an
obstruction is limited to that class. An integrated symbol is not the
integrand's canonical form.

\textbf{Deliverables.} Return either an explicit G, map, canonical form,
and exact equality check, or a proof excluding the declared ansatz
class. A failed search returns ``unresolved within budget.'' A dlog
decomposition alone is a candidate diagnostic, not a geometry
certificate.

\textbf{Validation.} Verify simplex/box normalization independently.
Check factorization and generalized cuts of the target. Rational
reconstruction must be followed by an exact polynomial identity after
denominators are cleared, with singular loci accounted for. Check
orientation signs and residues at infinity explicitly.

\textbf{Milestones.} Baseline: exact target and pole/residue audit.
Strong: one verified representation or one restricted no-go theorem.
Research extension: broaden the geometry class or multiplicity; explain
how the result changes understanding of gravity amplitudes rather than
treating a computational failure as a universal obstruction.

\href{https://github.com/rick-stevens-ai/Pure-Thought/blob/main/reviews/05-Positive-Geometry-Gravity.md}{Detailed
critique}.

\section{Synthesis 06. Non-perturbative S-matrix Bootstrap with Gravity}

\textbf{Mapping:}
\href{https://github.com/rick-stevens-ai/Pure-Thought/blob/main/PRDs/06-Nonperturbative-S-matrix-Bootstrap.md}{PRD
06}. Covered by canonical revision.

\subsection{06. S-matrix bootstrap with an explicit gravitational
infrared prescription}

\textbf{Question.} Can one certify an exclusion for a low-energy
scalar-scattering parameter using a crossing-symmetric partial-wave
relaxation whose infrared and truncation assumptions are explicit?

\textbf{Why it matters.} The method could identify restrictions shared
by many possible microscopic theories. That is useful when no single
theory is established. Its reliability depends on not mistaking a
limitation of the search formula, grid or cutoff for a limitation
imposed by physics.

\textbf{Scope and model.} Begin with identical massive scalars in four
dimensions and no gravity to validate the machinery. The gravitational
stage must first define an infrared-finite observable or a justified
regulator/pole treatment, including massless radiation. Keep one scalar
mass and one amplitude coefficient fixed; optimize a second coefficient
in a declared normalization.

This is complementary to Challenge 02: the focus is amplitude
feasibility and partial-wave constraints rather than deriving a
low-energy positivity sum rule. A nongravitational baseline does not
count as solving the gravitational stage.

\subsubsection{Mathematical formulation and claim boundary}

\href{https://github.com/rick-stevens-ai/Pure-Thought/blob/main/PRDs/06-Nonperturbative-S-matrix-Bootstrap.md}{Canonical
specification}.

\[A=16\pi\sum_{\ell\ge0}(2\ell+1)a_\ell(s)P_\ell(\cos\theta),\quad
\rho=\sqrt{1-4m^2/s},\quad S_\ell=1+2i\rho a_\ell.\]

\[|S_\ell|\le1\quad\Longleftrightarrow\quad
\operatorname{Im}a_\ell\ge\rho|a_\ell|^2\qquad(s>4m^2).\]

Equality requires an elastic sector. A certified outer relaxation can
exclude a physical amplitude only after the inclusion of every
admissible amplitude and all tail bounds are proved. Failure of an inner
ansatz excludes only that ansatz. The gravitational stage has a separate
infrared prerequisite.

\textbf{Deliverables.} Export the amplitude conventions, IR
prescription, finite optimization problem, mapping from the continuum
problem, and dual certificate with its positivity/error margins.
Feasible ansatz points are examples satisfying the verified constraints
only. If the IR construction or tail control fails, state which
conclusion remains unavailable.

\textbf{Validation.} Check free scattering and a specified analytic
nongravitational benchmark. Verify the unitarity disk algebra
independently. Test between grid points and beyond cutoffs using proved
bounds. Demonstrate that known admissible examples are not falsely
excluded. Increasing precision is a diagnostic; rigorous enclosures are
needed for certification.

\textbf{Milestones.} Baseline: certified finite nongravitational
relaxation. Strong: justify an IR treatment and certify one conditional
gravitational exclusion. Research extension: improved continuum bounds
or a constructive amplitude with a precisely delimited consistency
claim.

\href{https://github.com/rick-stevens-ai/Pure-Thought/blob/main/reviews/06-Nonperturbative-S-matrix-Bootstrap.md}{Detailed
critique}.

\section{Synthesis 07. Extremal Higher-Dimensional CFTs with Stress
Tensor}

\textbf{Mapping:}
\href{https://github.com/rick-stevens-ai/Pure-Thought/blob/main/PRDs/07-Extremal-CFTs-Stress-Tensor.md}{PRD
07}. Covered by canonical revision.

\subsection{07. Stress-tensor bootstrap: a sector-specific gap bound in
three dimensions}

\textbf{Question.} At fixed stress-tensor two- and three-point data, can
crossing exclude a proposed gap in the parity-even scalar sector of
T\ensuremath{\times}T in a unitary three-dimensional CFT?

\textbf{Why it matters.} The stress tensor is present in a broad class
of local theories, so restrictions derived from it can have wide reach.
Such restrictions help map which strongly interacting quantum theories
are possible. They can also constrain proposed boundary descriptions of
gravity, where the stress tensor has a special role.

\textbf{Scope and model.} Fix d=3, parity invariance, a unique conserved
stress tensor, a normalization for C\_T, and the remaining independent
TTT parameter. The initial observable is the first parity-even
nonidentity scalar appearing with nonzero coupling in
T\ensuremath{\times}T. It is not a bound on every operator in the
theory. State separate gaps for any other restricted spin/parity
sectors.

Use the full conserved spinning correlator. A scalar block toy
implementation is permitted for learning but is not the target
stress-tensor bootstrap.

\subsubsection{Mathematical formulation and claim boundary}

\href{https://github.com/rick-stevens-ai/Pure-Thought/blob/main/PRDs/07-Extremal-CFTs-Stress-Tensor.md}{Canonical
specification}.

\[V_{\mathrm{id}}+V_T(C_T,t_4)+
\sum_{\mathcal O}\lambda_{\mathcal O}^{\mathsf T}V_{\Delta,\ell,\pi}
\lambda_{\mathcal O}=0.\]

\[\alpha(V_{\mathrm{id}}+V_T)>0,\qquad
\alpha(V_{\Delta,\ell,\pi})\succeq0
\quad\text{for every allowed exchange}.\]

These inequalities contradict crossing and exclude the proposed
sector-specific gap. Tensor structures, conserved spinning blocks, Ward
identities and continuous exchange coverage are required. In three
dimensions a conserved symmetric-traceless spin-\(\ell\) primary has
\(\Delta=\ell+1\).

\textbf{Deliverables.} Export tensor/block conventions, conservation and
Ward identities, fixed input data, derivative basis, spin treatment, and
functional coefficients. Attach rigorous block-approximation and
high-spin control. If only the finite approximant is checked, label the
result accordingly.

\textbf{Validation.} Recover a selected published stress-tensor
benchmark using exactly its conventions and assumptions. Check the
allowed region rather than the excluded region. Verify matrix
positivity, not just individual scalar components. Free scalar/fermion
data provide normalization tests; numerical agreement alone is not a
proof of the continuum bound.

\textbf{Milestones.} Baseline: reproduce one finite benchmark and
publish the assumptions. Strong: turn one exclusion into a certified
result with omitted-sector control. Research extension: improve a bound,
add another sector, or obtain a conditional extremal spectrum.
Saturation does not construct a CFT.

\href{https://github.com/rick-stevens-ai/Pure-Thought/blob/main/reviews/07-Extremal-CFTs-Stress-Tensor.md}{Detailed
critique}.

\section{Synthesis 08. Swampland via Modularity \& Higher-Form
Symmetries}

\textbf{Mapping:}
\href{https://github.com/rick-stevens-ai/Pure-Thought/blob/main/PRDs/08-Swampland-Modularity-Symmetries.md}{PRD
08}. Covered by canonical revision.

\subsection{08. Exact modular-data and symmetry-anomaly consistency
tests}

\textbf{Question.} For supplied rational-CFT modular data and a
specified finite internal symmetry, which algebraic consistency tests
can be certified, and which failures obstruct a proposed gauging?

\textbf{Why it matters.} Exact consistency checks can detect false
candidate data before researchers invest in a full construction. They
also clarify which symmetries can be gauged and which require additional
structure, such as anomaly cancellation or inflow from another system.
This is useful both for studying quantum field theories and for
interpreting holographic gravity proposals.

\textbf{Scope and model.} Use a bosonic rational CFT with a declared
chiral algebra and finitely many characters. Input exact algebraic S,T
matrices, vacuum label, and a nonnegative integer multiplicity matrix M.
Start with the Ising model. An optional finite 0-form group G requires
separate defect/twisted-sector data and an anomaly cocycle; these cannot
be inferred from the untwisted partition function alone.

The result concerns the supplied CFT data and gauging problem. A general
bulk quantum-gravity compatibility verdict is outside the baseline.

\subsubsection{Mathematical formulation and claim boundary}

\href{https://github.com/rick-stevens-ai/Pure-Thought/blob/main/PRDs/08-Swampland-Modularity-Symmetries.md}{Canonical
specification}.

\[S^\dagger S=I,\quad S^2=C,\quad(ST)^3=C,\quad
MS=SM,\quad MT=TM,\quad M_{00}=1,\quad M_{ij}\in\mathbb Z_{\ge0}.\]

\[N_{ij}^{\ k}=\sum_a\frac{S_{ia}S_{ja}S_{ka}^*}{S_{0a}}
\in\mathbb Z_{\ge0},\qquad \omega\in Z^3(G,U(1)),\quad\delta\beta=\omega.\]

Fix the \(T\) convention and any assumed symmetry of \(S\). An explicit
cochain \(\beta\) trivializes the specified anomaly; a complete
cohomology calculation may prove no such cochain exists. This gauging
obstruction does not make an anomalous boundary theory inconsistent, and
untwisted modular data alone do not supply defect data.

\textbf{Deliverables.} Return pass/fail/unresolved for each named axiom,
with exact arithmetic traces. Include the number field and embeddings.
Even passing every modular-data test gives necessary algebraic
consistency only; full CFT sewing and realizability are separate
obligations.

\textbf{Validation.} Verify Ising S entries including
1/\ensuremath{\surd}2 in algebraic arithmetic. Reject altered modular
relations, noninteger multiplicities, and a failed cocycle identity.
Accept the existence of a boundary global symmetry rather than flagging
it as automatically inconsistent. Cross-check a small cyclic-group
cohomology example.

\textbf{Milestones.} Baseline: exact modular-data checker on Ising and
one second supplied model. Strong: one finite-group
twisted-sector/gauging analysis. Research extension: classify a
explicitly finite family of data or derive a conditional bulk
consequence using a stated holographic theorem.

\href{https://github.com/rick-stevens-ai/Pure-Thought/blob/main/reviews/08-Swampland-Modularity-Symmetries.md}{Detailed
critique}.

\section{Synthesis 09. Topological Band Theory Without Materials Data}

\textbf{Mapping:}
\href{https://github.com/rick-stevens-ai/Pure-Thought/blob/main/PRDs/09-Topological-Band-Theory.md}{PRD
09},
\href{https://github.com/rick-stevens-ai/Pure-Thought/blob/main/PRDs/15-Topological-Quantum-Chemistry.md}{PRD
15}. Covered by canonical revision.

\subsection{09. Certified Chern phases in a bounded tight-binding
family}

\textbf{Question.} Within a declared finite-range Hamiltonian family,
which parameter regions have a proven band gap and occupied-band Chern
number, and what minimality statements follow within that family?

\textbf{Why it matters.} Topological models explain why some properties
depend on global structure rather than microscopic details. A certified
phase map provides reliable examples for theory and simulation, and it
helps identify which ingredients a proposed design actually needs.
Proving a gap everywhere is important because a small missed gap closing
can invalidate a computed topological label.

\textbf{Scope and model.} Begin with the two-band square-lattice model
\(H_m(k)=\sin(k_x)\sigma_x+\sin(k_y)\sigma_y+[m+\cos(k_x)+\cos(k_y)]\sigma_z\),
k\ensuremath{\in}[\ensuremath{-}\ensuremath{\pi},\ensuremath{\pi}]\ensuremath{{}^{2}}.
Use one occupied band and a fixed orientation. Analyze rational
parameter intervals separated from m=\ensuremath{-}2,0,2.

For subsequent searches fix dimension, orbitals, filling, hopping
support, coefficient domain, and symmetry representation. A rational
grid is finite; a continuous coefficient box requires a different
completeness argument. ``All Hamiltonians for a space group'' is not an
enumerable specification.

\subsubsection{Mathematical formulation and claim boundary}

\href{https://github.com/rick-stevens-ai/Pure-Thought/blob/main/PRDs/09-Topological-Band-Theory.md}{Canonical
specification}.

\[H_m(k)=\sin k_x\,\sigma_x+\sin k_y\,\sigma_y+
(m+\cos k_x+\cos k_y)\sigma_z.\]

\[E_{r+1}(k)-E_r(k)\ge\delta>0,\qquad
C=\frac{1}{2\pi i}\int_{\mathrm{BZ}}
\operatorname{Tr}\!\left(P[\partial_{k_x}P,\partial_{k_y}P]\right)\,d^2k.\]

Prove the gap over the full claimed parameter and momentum domain. Fix
orientation and the occupied projector. An analytic argument or
validated integral enclosing one integer certifies \(C\); a momentum
mesh alone does not.

\textbf{Deliverables.} Return model coefficients, gap lower bounds,
invariant and sign convention, certified parameter coverage, and
uncovered boxes. A minimality claim must specify whether it minimizes
orbitals, hopping range, or coefficients and must exclude every smaller
object in that declared class.

\textbf{Validation.} Recover gap closings at \ensuremath{-}2,0,2 and
trivial/nontrivial intervals on either side. Reverse orientation and
verify C changes sign. A fine-grid lattice Chern calculation is a useful
cross-check, but must not substitute for the continuous gap proof. Check
a time-reversal-invariant occupied bundle has zero total Chern number.

\textbf{Milestones.} Baseline: complete phase diagram for H\_m away from
transition values. Strong: a certified atlas for one bounded
symmetry-compatible family. Research extension: new range/band tradeoffs
or minimality theorems, explicitly relative to the permitted family.

\href{https://github.com/rick-stevens-ai/Pure-Thought/blob/main/reviews/09-Topological-Band-Theory.md}{Detailed
critique}.

\subsection{15. Symmetry-indicator algebra and the limits of band
diagnosis}

\textbf{Question.} For one fixed symmetry setting, can compatibility
relations and atomic band representations be generated exactly, and what
topology can their quotient actually detect?

\textbf{Why it matters.} Symmetry labels are often much simpler to work
with than the full wavefunctions across momentum space. Knowing what
those labels do and do not prove can make topological diagnosis more
efficient and prevent false declarations that a system is ordinary. The
distinction is essential for turning a classification table into a
trustworthy scientific tool.

\textbf{Scope and model.} Choose one small symmetry setting, initially a
spinless two-dimensional inversion-symmetric lattice without time
reversal. Specify the space-group generators, orbital/site
representations, and high-symmetry momenta. Derive all input
representations from group operations, with external tables used only as
checks.

The result is an algebraic classification of symmetry data. A zero
indicator is not a certificate of a trivial occupied bundle.

\subsubsection{Mathematical formulation and claim boundary}

\href{https://github.com/rick-stevens-ai/Pure-Thought/blob/main/PRDs/15-Topological-Quantum-Chemistry.md}{Canonical
specification}.

\[A\subset B\subset\mathbb Z^r,\qquad X=B/A,\qquad
UA_{\mathrm{mat}}V=D,\quad U,V\text{ unimodular}.\]

\[b=A_{\mathrm{mat}}n,\quad n\in\mathbb Z^s
\qquad\text{versus}\qquad b=A_{\mathrm{mat}}n,\quad n\in\mathbb Z_{\ge0}^s.\]

Use a declared basis for the compatibility lattice \(B\) in the Smith
decomposition. Integer membership and nonnegative atomic symmetry-data
decomposition are distinct tests. Neither a zero indicator nor matching
atomic labels alone constructs localized Wannier functions for a given
Hamiltonian.

\textbf{Deliverables.} Export group conventions, irreps, compatibility
matrix, atomic generators, Smith decomposition, indicator mapping, and
witnesses for any integer/nonnegative decomposition. Report ``indicator
trivial; topology unresolved'' when symmetry data are insufficient. Use
Berry/Wilson-loop information when making a stronger claim about a
specific H.

\textbf{Validation.} Verify all representation identities and dimension
counts. Reconstruct the Smith identity U A\_mat V=D in the chosen B
basis, with U,V unimodular. Test known atomic data, incompatible data,
and models sharing symmetry labels but differing in an additional
invariant. Keep stable and fragile notions distinct.

\textbf{Milestones.} Baseline: exact quotient and decompositions for one
setting. Strong: connect the diagnosis to explicit Hamiltonians and
identify its kernel. Research extension: another symmetry group or a
rigorous fragile obstruction with a definition of allowed added bands.

\href{https://github.com/rick-stevens-ai/Pure-Thought/blob/main/reviews/15-Topological-Quantum-Chemistry.md}{Detailed
critique}.

\section{Synthesis 10. Flat Chern Bands with Provable Geometry}

\textbf{Mapping:}
\href{https://github.com/rick-stevens-ai/Pure-Thought/blob/main/PRDs/10-Flat-Chern-Bands.md}{PRD
10}. Covered by canonical revision.

\subsection{10. Flatness--locality--geometry tradeoffs for Chern bands}

\textbf{Question.} For fixed orbital count, hopping range, and gap
normalization, how small can a Chern band's bandwidth and geometric
nonuniformity be, with certified bounds?

\textbf{Why it matters.} Flat topological bands are useful settings for
studying collective quantum behavior driven by interactions. But there
are tradeoffs: restricting the model's spatial range changes which
flatness properties are possible. Quantifying those tradeoffs can show
whether a proposed improvement is attainable or requires a more
complicated model.

\textbf{Scope and model.} Choose a two-dimensional isolated rank-one
band, target C=1, and a compact coefficient family with range R and
fixed orbital positions. Initially use a specified two-band family
extending Challenge 09. Impose a positive gap floor and a norm bound to
prevent trivial energy rescaling.

Maintain two distinct tracks: finite-range approximate flatness, and
exact spectral flattening with generally non-finite-range hopping. Do
not require exact nonzero-Chern flatness with strictly finite-range
hopping under the no-go theorem's hypotheses.

\subsubsection{Mathematical formulation and claim boundary}

\href{https://github.com/rick-stevens-ai/Pure-Thought/blob/main/PRDs/10-Flat-Chern-Bands.md}{Canonical
specification}.

\[W=\max_k E_n(k)-\min_k E_n(k),\quad
\Delta=\min_k\min_{j\ne n}|E_j(k)-E_n(k)|,\quad\text{minimize }W/\Delta.\]

\[g_{ij}=\operatorname{Re}\langle\partial_i u|(1-P)|\partial_j u\rangle,
\quad \bar F=\frac{2\pi C}{A},\quad A=\operatorname{Area}(\mathrm{BZ}),\]

\[\operatorname{Var}(F)=\frac1A\int_{\mathrm{BZ}}(F-\bar F)^2d^2k,
\qquad D_g=\frac1A\int_{\mathrm{BZ}}(\operatorname{tr}g-|F|)d^2k.\]

Fix a gap floor, norm bound, orbital positions and coordinate metric.
Use the Chern convention of 09. Strictly finite-range approximate
flatness and generally long-range exact flattening are separate tasks;
these diagnostics alone do not prove an interacting fractional phase.

\textbf{Deliverables.} Export coefficient/domain constraints, W upper
and \ensuremath{\Delta} lower bounds, certified C, geometric enclosures,
and optimization gaps. An achieved objective gives an upper bound on the
minimum; a global lower bound needs a covering/relaxation proof. For
infinite-range flattening, give decay and truncation estimates.

\textbf{Validation.} Verify normalization
\ensuremath{\int}F=2\ensuremath{\pi}C and the pointwise metric
inequality in the chosen convention. Check that truncating a spectrally
flattened model generally reintroduces bandwidth. Include a trivial
exactly flat band as a negative topology test. Compare candidate
objectives using identical units, embedding, and norm constraints.

\textbf{Milestones.} Baseline: one certified near-flat Chern model and a
correct locality audit. Strong: global objective bounds on one compact
family or a rigorous truncation tradeoff. Research extension: optimal
geometry within a specified class; many-body FCI stability requires an
added interacting Hamiltonian and filling.

\href{https://github.com/rick-stevens-ai/Pure-Thought/blob/main/reviews/10-Flat-Chern-Bands.md}{Detailed
critique}.

\section{Synthesis 11. RBCE (Relativistic Band \& Crystal-Field
Engineering)}

\textbf{Mapping:}
\href{https://github.com/rick-stevens-ai/Pure-Thought/blob/main/supplemental/S01.md}{S01}.
Distinct objective preserved in supplement.

\subsection{S01. Crystal-field and spin--orbit anisotropy in a finite
ion model}

\textbf{Question.} For a fixed d-shell occupancy and point-group
representation, which anisotropy values are possible in a bounded
crystal-field/spin--orbit Hamiltonian family?

\textbf{Why it matters.} Directional magnetic response is a basic
ingredient in understanding magnetic materials and controlling spin
states. A finite ion model lets us separate the roles of local geometry
and spin--orbit coupling before tackling interactions among many ions.

\textbf{Scope and model.} Fix one ion model (start with
d\ensuremath{{}^{1}}), a chosen point group such as
C\ensuremath{{}_{4}}v, and a finite orbital-spin basis. Supply exact
matrices for L,S and the group action. Define
H(B)=H\_CF+\ensuremath{\lambda}L\ensuremath{\cdot}S+\ensuremath{\mu}\_B
B\ensuremath{\cdot}(L+2S), with crystal-field coefficients and
\ensuremath{\lambda} in a compact normalized box. Specify the
temperature or isolated low-energy multiplet before defining
susceptibility or the effective g tensor. For a Kramers doublet, extract
a g tensor from the projected magnetic moment and use its singular
values, which remove doublet-basis ambiguity. Require a certified
excitation gap above that doublet. Optimize a bounded ratio or
difference with a denominator bounded away from zero. Export matrices,
symmetry identities, interval eigenvalue bounds, and global objective
brackets.

\subsubsection{Mathematical formulation and claim boundary}

\href{https://github.com/rick-stevens-ai/Pure-Thought/blob/main/supplemental/S01.md}{Supplemental
specification}.

\[H(B)=H_{\mathrm{CF}}+\lambda L\cdot S+\mu_B B\cdot(L+2S),\]

\[P H(B)P=E_0I+\frac{\mu_B}{2}\sum_{a,b}B_a g_{ab}\sigma_b+O(|B|^2).\]

Use the isolated Kramers-doublet projector \(P\) and exact orbital-spin
conventions. The singular values of \(g\) remove doublet-basis
ambiguity. Certify the gap above the doublet and bound any denominator
in an anisotropy ratio on the compact parameter box. This finite ion
model does not establish bulk magnetism.

\textbf{Critique.} The synthesis's ``all d\^{}n configurations across
every point group'' has no interaction model, coefficient bounds, or
objective normalization. Symmetry permits tensor components but does not
determine their magnitudes; arbitrary crystal-field or spin--orbit
scaling can make an optimization meaningless. ``Rare-earth-like'' is a
motivation, not an observable. Many-electron d\^{}n ions also require
Coulomb/Hund parameters and a specified approximation. The revision
retains an anisotropy atlas but begins with a finite, normalized
spectral problem and separates a single-ion response from bulk
magnetocrystalline anisotropy.

\textbf{Validation and milestones.} Baseline: exact symmetry-allowed
Hamiltonian basis and one doublet g tensor. Strong: certified
parameter-box extrema with a doublet-gap condition. Extend to
d\ensuremath{{}^{2}} only after specifying interactions. Check the
isotropic limit, group covariance, and basis invariance of g singular
values.

Use the shared
\href{https://github.com/rick-stevens-ai/Pure-Thought/blob/main/reviews/EVIDENCE-STANDARD.md}{evidence
standard}. Unsuccessful search is unresolved unless a complete exclusion
proof is supplied. These are proposed baseline specifications, not
completed scientific results.

\section{Synthesis 12. Photonic/Phononic Crystals: Rigorous Bandgap
Optimization}

\textbf{Mapping:}
\href{https://github.com/rick-stevens-ai/Pure-Thought/blob/main/PRDs/11-Photonic-Topological-Crystals.md}{PRD
11},
\href{https://github.com/rick-stevens-ai/Pure-Thought/blob/main/supplemental/S02.md}{S02}.
Distinct objective preserved in supplement.

\subsection{11. Certified topological bands in an idealized Maxwell
medium}

\textbf{Question.} Can a specified lossless periodic electromagnetic
model support a rigorously isolated band group with nonzero Chern
number, and can its interface modes be validated?

\textbf{Why it matters.} Controlling where and in which direction waves
travel is a basic aim of optical engineering. Topology offers one
mechanism for robust interface behavior under appropriate conditions. A
certified Maxwell calculation would also expose whether an apparent
effect is caused by numerical artifacts or by assumptions a real
material cannot satisfy.

\textbf{Scope and model.} Choose a two-dimensional periodic, lossless,
frequency-independent constitutive model with Hermitian positive
\ensuremath{\epsilon}(r), \ensuremath{\mu}(r), specified boundary
conditions, and a mechanism breaking physical time reversal. Treat the
constitutive tensors as mathematical inputs. Real-frequency dispersion,
loss, and fabrication feasibility are extensions requiring additional
physical data.

Start by validating a reciprocal scalar dielectric benchmark with zero
total Chern number for a time-reversal-invariant band group; then
introduce an explicitly nonreciprocal tensor.

\subsubsection{Mathematical formulation and claim boundary}

\href{https://github.com/rick-stevens-ai/Pure-Thought/blob/main/PRDs/11-Photonic-Topological-Crystals.md}{Canonical
specification}.

\[(\nabla+ik)\times\mu^{-1}(r)(\nabla+ik)\times u
=\omega^2\epsilon(r)u,\qquad
\langle u,v\rangle_\epsilon=\int_{\mathrm{cell}}u^*\cdot\epsilon v\,dr.\]

Use positive Hermitian constitutive tensors, the transverse constraint,
units \(c=1\), and an isolated positive-frequency band group. Certify
its gap and projector invariant in the weighted metric. Nonzero total
Chern number requires breaking the relevant physical time-reversal
symmetry; scalar reciprocal benchmarks provide rejection controls.

\textbf{Deliverables.} Return ideal geometry and constitutive tensors,
solver discretization, eigenvalue enclosures, spurious-mode controls,
gap and topology certificates. Interface calculations must state both
media and termination. Report robustness only for an explicit
perturbation class and bound that preserves the required gap/symmetry.

\textbf{Validation.} Check homogeneous-medium dispersion, mesh/basis
convergence, Hermiticity, and transverse modes. Separate Galerkin error
from k-space interpolation error. Verify reciprocal controls cannot
acquire a net Chern number from numerical artifacts. Compare an
interface spectrum with the bulk-index difference without claiming
immunity to arbitrary disorder.

\textbf{Milestones.} Baseline: validated reciprocal eigenproblem.
Strong: one certified ideal nonreciprocal topological design. Research
extension: a realistic dispersive constitutive model or gap optimization
within a bounded geometry family. Fabrication drawings alone do not
certify material behavior.

\href{https://github.com/rick-stevens-ai/Pure-Thought/blob/main/reviews/11-Photonic-Topological-Crystals.md}{Detailed
critique}.

\subsection{S02. Bandgap optimization in a bounded wave-operator family}

\textbf{Question.} Within a fixed finite family of periodic media, what
lower bound on a complete bandgap and what upper bound on the best
possible gap can be certified?

\textbf{Why it matters.} Controlling which waves propagate supports
ideas for sound isolation, vibration management and optical filtering.
The same mathematical challenge arises across these applications: a
promising numerical design needs reliable evidence that its gap exists
everywhere and that any optimality claim covers the stated design space.

\textbf{Scope and model.} Choose either a scalar acoustic operator or a
full Maxwell operator, never an unspecified mixture. Initially choose a
periodic scalar elliptic operator
\ensuremath{-}\ensuremath{\nabla}\ensuremath{\cdot}a(r)\ensuremath{\nabla}
with 0<a\_min\ensuremath{\le}a(r)\ensuremath{\le}a\_max, a fixed unit
cell and a finite geometry parameter box. For adjacent Bloch eigenvalues
define the complete gap g=min\_k
\ensuremath{\lambda}\_\{n+1\}(k)\ensuremath{-}max\_k
\ensuremath{\lambda}\_n(k), with a declared normalized objective. A
grid-local separation is not a complete gap. Use validated eigenvalue
bounds plus k-space variation estimates to cover the Brillouin zone.
Branch-and-bound over parameters yields a candidate lower bound and an
optimization upper bound, keeping the inequality direction explicit.

\subsubsection{Mathematical formulation and claim boundary}

\href{https://github.com/rick-stevens-ai/Pure-Thought/blob/main/supplemental/S02.md}{Supplemental
specification}.

\[\lambda_- =\max_k\lambda_n(k),\qquad
\lambda_+=\min_k\lambda_{n+1}(k),\qquad
g=\lambda_+-\lambda_-. \]

A possible normalized objective is \(G=2g/(\lambda_++\lambda_-)\) when
the denominator is positive; frequency-gap objectives must instead use
\(\omega_j=\sqrt{\lambda_j}\) consistently.

Use a fixed self-adjoint wave model, bounded coefficients and geometry
class, with a positive denominator. A complete gap needs
\(\lambda_+>\lambda_-\) across the entire Brillouin zone and all
required polarizations. A certified candidate lower bound does not
establish global optimality; that needs coverage of the full design
class.

\textbf{Critique.} The old description mixes photonic and phononic
operators, index contrast and generic gap optimization without fixing
constitutive laws, geometry regularity, or constraints. ``Lowest
possible contrast'' and ``largest possible gap'' are two objectives
needing an ordering or Pareto definition. A converged plane-wave plot
does not establish a complete gap, and a locally optimized geometry does
not prove global optimality. This is distinct from PRD 11's topological
invariant: a wide ordinary bandgap can be topologically trivial.

\textbf{Validation and milestones.} Baseline: one certified complete gap
for a supplied geometry. Strong: bound the optimum to a declared
tolerance over the entire compact parameter box, or report unresolved
boxes. Validate a homogeneous medium and a simple one-dimensional
layered case before multidimensional optimization.

Use the shared
\href{https://github.com/rick-stevens-ai/Pure-Thought/blob/main/reviews/EVIDENCE-STANDARD.md}{evidence
standard}. Unsuccessful search is unresolved unless a complete exclusion
proof is supplied. These are proposed baseline specifications, not
completed scientific results.

\section{Synthesis 13. Universal Bounds for Effective Properties of
Composites}

\textbf{Mapping:}
\href{https://github.com/rick-stevens-ai/Pure-Thought/blob/main/supplemental/S03.md}{S03}.
Distinct objective preserved in supplement.

\subsection{S03. Effective conductivity bounds for two-phase composites}

\textbf{Question.} For fixed positive phase conductivities and volume
fraction, can rigorous effective-conductivity bounds and matching
constructions be obtained in a specified microstructure class?

\textbf{Why it matters.} Composite design seeks useful bulk properties
from available ingredients. Bounds reveal whether a proposed arrangement
performs well, whether a computation is accurate, and how much
improvement remains possible. Simple layered structures also provide
exact reference cases for checking more complicated methods.

\textbf{Scope and model.} Choose scalar conductivities
a\ensuremath{{}_{1}},a\ensuremath{{}_{2}}>0, periodic unit cell Y, phase
fraction f, and a unit direction e. Define e\ensuremath{{}^{T}}A\_eff
e=min\_\{periodic
u\}\ensuremath{\langle}a|e+\ensuremath{\nabla}u|\ensuremath{{}^{2}}\ensuremath{\rangle}.
Use an admissible potential trial field for an upper bound and a
divergence-free flux variational principle for a lower bound.
Distinguish bounds valid for every arrangement from optimization over a
finite geometry family. Fix isotropy assumptions before comparing
isotropic bounds. Cross-property elasticity/thermal results are separate
extensions with both constitutive laws specified.

\subsubsection{Mathematical formulation and claim boundary}

\href{https://github.com/rick-stevens-ai/Pure-Thought/blob/main/supplemental/S03.md}{Supplemental
specification}.

\[e\cdot A^*e=\min_{\phi\ \mathrm{periodic}}
\langle a(x)|e+\nabla\phi|^2\rangle,\qquad a(x)\in\{a_1,a_2\},\]

\[\left(\frac f{a_1}+\frac{1-f}{a_2}\right)^{-1}|e|^2
\le e\cdot A^*e\le(fa_1+(1-f)a_2)|e|^2.\]

Assume positive scalar phases, fixed volume fraction \(f\), and a
normalized cell average. Potential trials give upper bounds and the
corresponding flux principle gives lower bounds. A finite microstructure
list does not characterize an unrestricted G-closure.

\textbf{Critique.} The original ``universal bounds'' problem spans
conductivity, permittivity and elasticity, including tensor and
frequency effects with different variational principles. ``Beyond
Hashin--Shtrikman'' is not a measurable novelty requirement until the
symmetry class and available information are fixed. A finite set of
candidate microstructures cannot characterize the unrestricted
G-closure. The revision makes one scalar static problem precise and
makes the bound direction and construction obligation explicit.

\textbf{Validation and milestones.} Baseline: verify arithmetic and
harmonic directional bounds with layered composites. Strong: matching
bounds for one declared geometry class or a certified gap to the
unrestricted bound. Check volume fraction and ellipticity exactly;
report discretization and geometry errors separately.

Use the shared
\href{https://github.com/rick-stevens-ai/Pure-Thought/blob/main/reviews/EVIDENCE-STANDARD.md}{evidence
standard}. Unsuccessful search is unresolved unless a complete exclusion
proof is supplied. These are proposed baseline specifications, not
completed scientific results.

\section{Synthesis 14. Topological Mechanics: Maxwell Frames \&
Programmable Metamaterials}

\textbf{Mapping:}
\href{https://github.com/rick-stevens-ai/Pure-Thought/blob/main/PRDs/12-Topological-Mechanical-Metamaterials.md}{PRD
12}. Covered by canonical revision.

\subsection{12. Rigidity and boundary modes of periodic Maxwell frames}

\textbf{Question.} For a geometrically specified isostatic frame and
boundary termination, can the zero-mode count and localization be proved
from its compatibility matrix?

\textbf{Why it matters.} Geometry-controlled softness could help explain
and design structures with directional compliance, boundary motion or
localized mechanical response. The scientific value is understanding
which behaviors follow from a robust index and which are accidental
features of one geometry. That distinction matters before translating an
ideal spring model into a physical structure.

\textbf{Scope and model.} Use unstressed central-force springs with
exact rational or algebraic site coordinates, positive spring constants,
and explicit periodic bond offsets. Fix the unit-cell gauge, dimension,
masses, and boundary termination. Begin with a small frame and a known
periodic Maxwell lattice.

Connectivity alone is insufficient: bond directions determine the
compatibility matrix. Distinguish rigid-body motions, floppy modes, and
states of self stress. Claims concern linearized ideal mechanics.

\subsubsection{Mathematical formulation and claim boundary}

\href{https://github.com/rick-stevens-ai/Pure-Thought/blob/main/PRDs/12-Topological-Mechanical-Metamaterials.md}{Canonical
specification}.

\[Q=C^\dagger,\qquad n_0-n_{ss}=dN-N_b,\qquad
K=C^\dagger\operatorname{diag}(k_b)C,\quad D=M^{-1/2}KM^{-1/2}.\]

\[\nu_\gamma=\frac1{2\pi i}\oint_\gamma d\log\det C(k),
\qquad \det C(k)\ne0\text{ on }\gamma.\]

The first count includes rigid-body motions and self stress. Supply
actual bond directions, unit-cell gauge and termination. Acoustic zeros
and the local boundary count must be included before identifying edge
modes from winding.

\textbf{Deliverables.} Export positions, bonds, C, exact rank/kernel
calculations, self-stress basis, winding certificate, and a
boundary-specific mode count. For localization, provide a
transfer-matrix/root-modulus bound separated from the unit circle.
Specify perturbations preserving the nonvanishing condition and
positivity of spring constants.

\textbf{Validation.} Check one elementary underconstrained and one
overconstrained frame against the index equation. Verify K is positive
semidefinite and translations are zero modes where appropriate. Compare
two terminations of the same bulk. Reject winding calculations crossing
a zero of det C.

\textbf{Milestones.} Baseline: exact Maxwell--Calladine counts on finite
frames. Strong: one periodic frame with certified winding and
termination-dependent edge localization. Research extension: a bounded
family with robust mode counts. Elastic beams, prestress, nonlinear
stability, and manufacturing tolerance require separate models.

\href{https://github.com/rick-stevens-ai/Pure-Thought/blob/main/reviews/12-Topological-Mechanical-Metamaterials.md}{Detailed
critique}.

\section{Synthesis 15. Real-Space Topological Invariants for Disordered
Media}

\textbf{Mapping:}
\href{https://github.com/rick-stevens-ai/Pure-Thought/blob/main/PRDs/09-Topological-Band-Theory.md}{PRD
09},
\href{https://github.com/rick-stevens-ai/Pure-Thought/blob/main/supplemental/S04.md}{S04}.
Related PRD; supplement preserves distinct objective.

\subsection{09. Certified Chern phases in a bounded tight-binding
family}

\textbf{Question.} Within a declared finite-range Hamiltonian family,
which parameter regions have a proven band gap and occupied-band Chern
number, and what minimality statements follow within that family?

\textbf{Why it matters.} Topological models explain why some properties
depend on global structure rather than microscopic details. A certified
phase map provides reliable examples for theory and simulation, and it
helps identify which ingredients a proposed design actually needs.
Proving a gap everywhere is important because a small missed gap closing
can invalidate a computed topological label.

\textbf{Scope and model.} Begin with the two-band square-lattice model
\(H_m(k)=\sin(k_x)\sigma_x+\sin(k_y)\sigma_y+[m+\cos(k_x)+\cos(k_y)]\sigma_z\),
k\ensuremath{\in}[\ensuremath{-}\ensuremath{\pi},\ensuremath{\pi}]\ensuremath{{}^{2}}.
Use one occupied band and a fixed orientation. Analyze rational
parameter intervals separated from m=\ensuremath{-}2,0,2.

For subsequent searches fix dimension, orbitals, filling, hopping
support, coefficient domain, and symmetry representation. A rational
grid is finite; a continuous coefficient box requires a different
completeness argument. ``All Hamiltonians for a space group'' is not an
enumerable specification.

\subsubsection{Mathematical formulation and claim boundary}

\href{https://github.com/rick-stevens-ai/Pure-Thought/blob/main/PRDs/09-Topological-Band-Theory.md}{Canonical
specification}.

\[H_m(k)=\sin k_x\,\sigma_x+\sin k_y\,\sigma_y+
(m+\cos k_x+\cos k_y)\sigma_z.\]

\[E_{r+1}(k)-E_r(k)\ge\delta>0,\qquad
C=\frac{1}{2\pi i}\int_{\mathrm{BZ}}
\operatorname{Tr}\!\left(P[\partial_{k_x}P,\partial_{k_y}P]\right)\,d^2k.\]

Prove the gap over the full claimed parameter and momentum domain. Fix
orientation and the occupied projector. An analytic argument or
validated integral enclosing one integer certifies \(C\); a momentum
mesh alone does not.

\textbf{Deliverables.} Return model coefficients, gap lower bounds,
invariant and sign convention, certified parameter coverage, and
uncovered boxes. A minimality claim must specify whether it minimizes
orbitals, hopping range, or coefficients and must exclude every smaller
object in that declared class.

\textbf{Validation.} Recover gap closings at \ensuremath{-}2,0,2 and
trivial/nontrivial intervals on either side. Reverse orientation and
verify C changes sign. A fine-grid lattice Chern calculation is a useful
cross-check, but must not substitute for the continuous gap proof. Check
a time-reversal-invariant occupied bundle has zero total Chern number.

\textbf{Milestones.} Baseline: complete phase diagram for H\_m away from
transition values. Strong: a certified atlas for one bounded
symmetry-compatible family. Research extension: new range/band tradeoffs
or minimality theorems, explicitly relative to the permitted family.

\href{https://github.com/rick-stevens-ai/Pure-Thought/blob/main/reviews/09-Topological-Band-Theory.md}{Detailed
critique}.

\subsection{S04. Finite-volume certification of a disordered topological
index}

\textbf{Question.} For a local disordered lattice Hamiltonian with a
specified bulk spectral-gap hypothesis, when does a finite-volume index
provably equal its infinite-volume counterpart?

\textbf{Why it matters.} Topological behavior is interesting partly
because it can survive certain imperfections. Establishing that
robustness requires methods that work without a perfect crystal and
distinguish a finite-sample observation from a bulk conclusion. A
rigorous bridge between the two makes numerical evidence more
informative.

\textbf{Scope and model.} Fix a finite-range two-dimensional
tight-binding Hamiltonian on Z\ensuremath{{}^{2}}, bounded disorder,
Fermi energy E\_F and a bulk spectral gap. Initially use bounded
perturbations small enough to preserve a known clean gap. Choose one
index construction, such as a spectral localizer, and one theorem
relating finite and infinite volume. Define position operators,
truncation radius and boundary convention; compute the finite Hermitian
signature or index with interval spectral bounds. Verify every locality,
commutator, radius and scale inequality in the selected theorem.
Mobility-gap cases require a different theorem and are an extension.

\subsubsection{Mathematical formulation and claim boundary}

\href{https://github.com/rick-stevens-ai/Pure-Thought/blob/main/supplemental/S04.md}{Supplemental
specification}.

\[\operatorname{dist}(E_F,\operatorname{spec}H_0)\ge g>0,
\qquad\|V\|<g\quad\Longrightarrow\quad
\operatorname{dist}(E_F,\operatorname{spec}(H_0+V))\ge g-\|V\|.\]

This bounded self-adjoint perturbation estimate supplies a spectral-gap
hypothesis. Equality of a chosen finite-volume index with the bulk index
still requires a specific localizer or other index theorem with all
radius, locality, commutator and scale conditions checked. Mobility-gap
claims need a different theorem.

\textbf{Critique.} The synthesis lists Bott indices, noncommutative
Chern numbers and localizers as if they shared one universal finite-size
error bound. They have different hypotheses and finite-volume behavior.
An integer from a finite matrix is not automatically the bulk invariant,
and one disorder realization does not establish robustness for all
disorder. The revised task certifies the finite-to-infinite transfer
under a stated gap/locality theorem; it is distinct from the periodic
model classification in PRD 09.

\textbf{Validation and milestones.} Baseline: recover a clean model's
known index with a validated finite-to-bulk theorem. Strong: certify a
full bounded disorder class using gap preservation. Include trivial
controls and a gap-closing case that the checker declines to certify.

Use the shared
\href{https://github.com/rick-stevens-ai/Pure-Thought/blob/main/reviews/EVIDENCE-STANDARD.md}{evidence
standard}. Unsuccessful search is unresolved unless a complete exclusion
proof is supplied. These are proposed baseline specifications, not
completed scientific results.

\section{Synthesis 16. N-Representability \& 2-RDM Variational
Chemistry}

\textbf{Mapping:}
\href{https://github.com/rick-stevens-ai/Pure-Thought/blob/main/PRDs/16-N-Representability-2RDM.md}{PRD
16}. Covered by canonical revision.

\subsection{16. Certified energy brackets from reduced-density-matrix
relaxations}

\textbf{Question.} For a specified finite fermionic Hamiltonian, how
tight an independently certified ground-state energy interval can P,Q,G
and selected higher-order RDM constraints produce?

\textbf{Why it matters.} Reliable energy calculations underpin the
theoretical study of bonding and interacting electrons. A bracket tells
us how much uncertainty remains, rather than merely reporting a
convincing decimal number. The approach may extract useful information
without storing a complete many-electron wavefunction, although general
representability remains computationally difficult.

\textbf{Scope and model.} Fix N\ensuremath{\ge}2 fermions in M spin
orbitals with an exact rational/algebraic, number-conserving one- and
two-body Hamiltonian. Start with a small Hubbard-type system where exact
diagonalization in the same N-particle sector is feasible. Molecular
integrals are a later input with an explicit basis and numerical error
enclosure.

Use ensemble N-representability. The task is to certify bounds within a
finite orbital model, not solve general electronic structure in
polynomial time.

\subsubsection{Mathematical formulation and claim boundary}

\href{https://github.com/rick-stevens-ai/Pure-Thought/blob/main/PRDs/16-N-Representability-2RDM.md}{Canonical
specification}.

\[\gamma_{ij}=\langle a_i^\dagger a_j\rangle,\qquad
\Gamma_{ij,kl}=\langle a_i^\dagger a_j^\dagger a_l a_k\rangle,\]

\[\operatorname{Tr}\gamma=N,\quad
\sum_{ij}\Gamma_{ij,ij}=N(N-1),\quad
\sum_j\Gamma_{ij,kj}=(N-1)\gamma_{ik}.\]

\[L_{\mathrm{dual}}\le E_{\mathrm{relax}}\le E_0\le
\langle\Psi|H|\Psi\rangle=U,\qquad \langle\Psi|\Psi\rangle=1.\]

Derive two-particle, two-hole and particle-hole moment matrices from
anticommutation in this convention. Necessary PSD constraints give an
outer relaxation. Certify dual feasibility for \(L\) and a physical
\(N\)-fermion state for \(U\); a relaxed RDM need not represent any
state.

\textbf{Deliverables.} Export the Hamiltonian and indexing convention,
moment constraints, certified dual objective L, state witness for U, and
gap U\ensuremath{-}L. Separate solver error, relaxation error, and
finite-basis/model error. Report any property bounds via separately
optimized observables; a relaxed \ensuremath{\Gamma} need not be
physical.

\textbf{Validation.} Compare with exact diagonalization in identical
conventions. Check contractions using explicit Slater determinants and
correlated states. Adding valid necessary constraints must not lower the
exact relaxed minimum. Verify dual PSD matrices and equality residual
corrections rigorously, rather than interpreting the raw solver
objective as a bound.

\textbf{Milestones.} Baseline: valid L,U brackets for small fixed
models. Strong: quantify improvement from G or T constraints across a
specified instance set. Research extension: a new valid constraint with
demonstrated tightening, or a theorem on a tractable family; novelty
needs comparison with existing hierarchies.

\href{https://github.com/rick-stevens-ai/Pure-Thought/blob/main/reviews/16-N-Representability-2RDM.md}{Detailed
critique}.

\section{Synthesis 17. Non-empirical Density Functional Theory via
Convex Analysis}

\textbf{Mapping:}
\href{https://github.com/rick-stevens-ai/Pure-Thought/blob/main/supplemental/S05.md}{S05}.
Distinct objective preserved in supplement.

\subsection{S05. Finite-basis density functionals and exact convex
constraints}

\textbf{Question.} For a fixed finite lattice Hamiltonian, can the
ensemble constrained-search density functional be bracketed and its
convexity/duality verified?

\textbf{Why it matters.} Density-based methods reduce a difficult
many-electron description to a much smaller set of variables. Exact
finite examples reveal which constraints an approximation should respect
and provide benchmarks for evaluating proposed energy formulas.

\textbf{Scope and model.} Fix a finite N-particle Hilbert space, kinetic
T, interaction W, and commuting site-density operators \(\widehat n_i\).
Define
\(F(n)=\min_{\rho\succeq0,\operatorname{Tr}\rho=1,\operatorname{Tr}(\rho\widehat n_i)=n_i}\operatorname{Tr}\rho(T+W)\)
on ensemble-representable densities. Derive the dual
\(F(n)=\sup_v[E(v)-v\cdot n]\) with
\(E(v)=\min_{\rho\succeq0,\operatorname{Tr}\rho=1}\operatorname{Tr}\rho(T+W+\sum_i v_i\widehat n_i)\),
under the finite-dimensional convex formulation. Use exact states for
primal upper bounds and potentials with certified E(v) lower bounds for
dual lower bounds. If defining an exchange-correlation remainder,
specify T\_s and Hartree conventions separately.

\subsubsection{Mathematical formulation and claim boundary}

\href{https://github.com/rick-stevens-ai/Pure-Thought/blob/main/supplemental/S05.md}{Supplemental
specification}.

\[F(n)=\min_{\substack{\rho\succeq0,\ \operatorname{Tr}\rho=1\\
\operatorname{Tr}(\rho\widehat n_i)=n_i}}
\operatorname{Tr}\rho(T+W),\]

\[E(v)=\min_{\rho\succeq0,\ \operatorname{Tr}\rho=1}
\operatorname{Tr}\rho\!\left(T+W+\sum_i v_i\widehat n_i\right),
\qquad F(n)=\sup_v[E(v)-v\cdot n].\]

Work in the fixed finite \(N\)-particle Hilbert space and on
ensemble-representable densities. The supremum need not be attained by a
finite potential at boundary densities. Feasible density matrices give
upper bounds; certified energy lower bounds at chosen potentials give
lower bounds. Convexity of \(F\) does not imply convexity of an
exchange-correlation remainder.

\textbf{Critique.} ``The cone of exchange-correlation functionals
satisfying all exact constraints'' is not justified: normalization and
affine conditions need not define a cone, and convexity of the ensemble
universal functional does not imply convexity of E\_xc after subtracting
other functionals. Constraint satisfaction alone does not identify the
exact physical functional. The broad continuum task also leaves
admissible density spaces and spin scaling unspecified. The revision
first tests exact convex duality in a finite model, preserving
nonempirical functional development as an extension.

\textbf{Validation and milestones.} Baseline: exact functional values
and dual brackets on a two-site model. Strong: a certified piecewise
characterization or a new valid inequality for a larger finite model.
Verify particle-number bounds, convexity inequalities, and agreement
with direct constrained search.

Use the shared
\href{https://github.com/rick-stevens-ai/Pure-Thought/blob/main/reviews/EVIDENCE-STANDARD.md}{evidence
standard}. Unsuccessful search is unresolved unless a complete exclusion
proof is supplied. These are proposed baseline specifications, not
completed scientific results.

\section{Synthesis 18. Strictly-Correlated Electrons as Multi-Marginal
Optimal Transport}

\textbf{Mapping:}
\href{https://github.com/rick-stevens-ai/Pure-Thought/blob/main/PRDs/18-Optimal-Transport-Chemistry.md}{PRD
18}. Covered by canonical revision.

\subsection{18. Certified optimal-transport bounds for strictly
correlated electrons}

\textbf{Question.} For a supplied normalized electron density, can a
discretized Coulomb multi-marginal transport problem yield verified
bounds, with continuum error separated from optimization error?

\textbf{Why it matters.} Strongly correlated electrons are difficult to
approximate because their motions are strongly linked. This limiting
problem gives a mathematically structured reference for studying that
linkage and for developing density-functional ideas. Verified bounds
help distinguish a true property of the continuum model from an artifact
of a spatial grid or optimization tolerance.

\textbf{Scope and model.} Take an electron density
\ensuremath{\rho}\ensuremath{\ge}0 with
\ensuremath{\int}\ensuremath{\rho}=N and marginal
\ensuremath{\mu}=\ensuremath{\rho}/N. Begin with N=2 and a compactly
supported mathematical density or a finite rational discrete model.
Specify treatment of the diagonal Coulomb singularity. A
continuous-density result additionally needs domain-tail and
discretization estimates.

The main target is the strictly correlated interaction functional.
Wasserstein comparison of densities is a separate optional metric task;
it is not a minimum-energy chemical reaction path.

\subsubsection{Mathematical formulation and claim boundary}

\href{https://github.com/rick-stevens-ai/Pure-Thought/blob/main/PRDs/18-Optimal-Transport-Chemistry.md}{Canonical
specification}.

\[\int\rho=N,\qquad \mu=\rho/N,\qquad
c(r_1,\ldots,r_N)=\sum_{i<j}|r_i-r_j|^{-1},\]

\[V_{\mathrm{SCE}}[\rho]=\inf_{\Gamma\in\Pi(\mu,\ldots,\mu)}\int c\,d\Gamma,
\qquad \sum_i u(r_i)\le c(r_1,\ldots,r_N),\]

\[\underbrace{\int u\rho}_{L}\ \le\ V_{\mathrm{SCE}}[\rho]\ \le\
\underbrace{\int c\,d\Gamma_{\mathrm{feasible}}}_{U}.\]

The coupling is a probability measure; it may be symmetrized and is not
antisymmetric. This is a Coulomb interaction functional, not exact
exchange. For a discrete problem, the same inequalities apply only to
the discrete optimum until continuum error is bounded.

\textbf{Deliverables.} Return marginal conventions, cost and singularity
policy, primal plan, dual potentials, rigorously checked constraints,
and L\ensuremath{\le}V\ensuremath{\le}U. State whether V is a finite
discrete optimum or the continuum functional. Sinkhorn candidates must
be converted to feasible bounds for the unregularized problem, with
entropic bias controlled.

\textbf{Validation.} Verify a small rational-cost transport LP exactly
and a one-dimensional quadratic-cost quantile example as a separate
solver test. For Coulomb densities, use a case with a known co-motion
construction under its actual hypotheses. Reject comparing normalized
densities with unequal total mass unless normalization or an
unbalanced-OT model is explicitly chosen.

\textbf{Milestones.} Baseline: certified two-marginal discrete bounds.
Strong: one continuum enclosure with singularity/tail control. Research
extension: N\textgreater2 with symmetry reduction and validated bounds,
or a demonstrably improved relaxation. Do not promise deterministic
Monge maps for arbitrary multimarginal inputs.

\href{https://github.com/rick-stevens-ai/Pure-Thought/blob/main/reviews/18-Optimal-Transport-Chemistry.md}{Detailed
critique}.

\section{Synthesis 19. Complete Isomer Enumeration with Proofs}

\textbf{Mapping:}
\href{https://github.com/rick-stevens-ai/Pure-Thought/blob/main/PRDs/17-Isomer-Enumeration-Molecular-Graphs.md}{PRD
17}. Covered by canonical revision.

\subsection{17. Complete enumeration of a precisely defined
molecular-graph class}

\textbf{Question.} For a fixed formula and explicit valence grammar, can
every connected constitutional graph be generated exactly once, with a
checkable completeness argument?

\textbf{Why it matters.} Complete enumeration turns an intuitive picture
of ``possible structures'' into a checkable mathematical set. Such sets
are useful for testing structure-generation methods and exploring how
bonding constraints shape chemical possibilities. The guarantee is
especially valuable when a search is used to support a claim that no
alternative structure exists.

\textbf{Scope and model.} Initially allow neutral closed-shell C,H,O
graphs with bond orders 1,2,3, no self-bonds, and valences 4,1,2.
Hydrogens may be implicit if reconstructed uniquely from residual
valence. Specify whether different Kekulé graphs are distinct; the
baseline uses literal bond-order graphs and no aromatic equivalence.

Constitutional graph validity is not chemical stability or
synthesizability. Stereoisomers, conformers, ions, radicals, isotopes,
and 3D geometries require added definitions and are not baseline
outputs.

\subsubsection{Mathematical formulation and claim boundary}

\href{https://github.com/rick-stevens-ai/Pure-Thought/blob/main/PRDs/17-Isomer-Enumeration-Molecular-Graphs.md}{Canonical
specification}.

Let \(\mathcal G\) contain exactly the graphs satisfying the fixed atom,
bond, valence and connectivity grammar. Then

\[\text{catalog}=\mathcal G/\!\cong.\]

Isomorphisms preserve the declared atom and bond labels. Every output
must satisfy the grammar, no two representatives may be equivalent, and
every allowed input must map to a representative. Canonical labeling
proves deduplication only; completeness needs a coverage proof for
generation and pruning. Stereochemistry requires a separately declared
equivalence model.

\textbf{Deliverables.} Export canonical graphs, reconstructed hydrogens,
counts, labeling conventions, and generation/coverage proof. Provide a
validator independent of the optimized generator. If a resource limit
stops enumeration, output a partial catalog and coverage status rather
than ``all isomers.''

\textbf{Validation.} Recover alkane constitutional counts for
C\ensuremath{{}_{4}}H\ensuremath{{}_{10}} and
C\ensuremath{{}_{5}}H\ensuremath{{}_{12}}, namely 2 and 3. Cross-check
all tiny instances with brute force. Random relabelings must preserve
canonical codes. Verify each graph's formula and weighted valence;
reject disconnected mixtures.

\textbf{Milestones.} Baseline: complete small-formula catalog with
independent counts. Strong: complete
C\ensuremath{{}_{6}}H\ensuremath{{}_{12}}O under the stated grammar.
Research extension: stereochemical orbit enumeration with meso cases and
explicit stereocenter/double-bond conventions; keep conformational
continua separate.

\href{https://github.com/rick-stevens-ai/Pure-Thought/blob/main/reviews/17-Isomer-Enumeration-Molecular-Graphs.md}{Detailed
critique}.

\section{Synthesis 20. Inverse Statistical Mechanics: Potentials with
Provable Ground States}

\textbf{Mapping:}
\href{https://github.com/rick-stevens-ai/Pure-Thought/blob/main/supplemental/S06.md}{S06}.
Distinct objective preserved in supplement.

\subsection{S06. Inverse pair-potential design over an explicit
competitor class}

\textbf{Question.} Can a parameterized pair potential make a target
periodic structure an energy minimizer among a precisely specified set
of competitors over a density interval?

\textbf{Why it matters.} Interaction design is a way to explore
self-assembly: simple local rules may favor organized structures.
Reliable comparisons expose whether apparent success comes from a useful
mechanism or from overlooking a competing arrangement or a long-range
contribution.

\textbf{Scope and model.} Fix dimension, density interval, and
v\_\ensuremath{\theta}(r)=\ensuremath{\Sigma}\_j\ensuremath{\theta}\_j\ensuremath{\phi}\_j(r)
with bounded coefficients and convergent lattice sums. Define
E\_\ensuremath{\theta}(\ensuremath{\Lambda})=1/2
\ensuremath{\Sigma}\_\{x\ensuremath{\in}\ensuremath{\Lambda}\ensuremath{\setminus}\{0\}\}v\_\ensuremath{\theta}(|x|).
Begin with a finite list of normalized lattice competitors and use
interval sums plus tail bounds to certify
E\_target\ensuremath{\le}E\_competitor\ensuremath{-}\ensuremath{\delta}
uniformly on a density interval. Continuous competitors require a
compact parameterization and global coverage. State whether particle
displacements or variable cell size are allowed. General configuration
ground states require additional universal lower bounds beyond lattice
comparisons.

\subsubsection{Mathematical formulation and claim boundary}

\href{https://github.com/rick-stevens-ai/Pure-Thought/blob/main/supplemental/S06.md}{Supplemental
specification}.

\[v_\theta(r)=\sum_j\theta_j\phi_j(r),\qquad
E_\theta(\Lambda)=\frac12\sum_{x\in\Lambda\setminus\{0\}}v_\theta(|x|),\]

\[E_\theta(\Lambda_{\mathrm{target}})\le
E_\theta(\Lambda_{\mathrm{competitor}})-\delta.\]

Require convergent sums and certified tails, uniformly on the declared
density interval and competitor set. The conclusion is best among those
competitors. Unrestricted configurations, infinitesimal deformations and
finite-temperature free energies require separate arguments.

\textbf{Critique.} The original jumps from optimizing lattice energies
to provable crystals over both densities and temperatures. Beating
selected lattices does not exclude aperiodic competitors or
infinitesimal deformations, and a zero-temperature energy gap is not a
finite-temperature free-energy theorem. A strictly uniform energy gap
against all arbitrarily close deformations may be impossible under
continuity. The rewrite gives a restricted inverse-design result whose
proof scope is honest; unrestricted crystallization and thermal
stability remain separate research extensions.

\textbf{Validation and milestones.} Baseline: exact or interval-verified
ordering for a finite competitor list. Strong: uniform
density/parameter-region certificate. Validate a direct long-cutoff sum
against the certified tail bound; label the result ``best among these
competitors'' unless a universal proof is supplied.

Use the shared
\href{https://github.com/rick-stevens-ai/Pure-Thought/blob/main/reviews/EVIDENCE-STANDARD.md}{evidence
standard}. Unsuccessful search is unresolved unless a complete exclusion
proof is supplied. These are proposed baseline specifications, not
completed scientific results.

\section{Synthesis 21. LDPC \& Hypergraph-Product Quantum Error
Correcting Codes}

\textbf{Mapping:}
\href{https://github.com/rick-stevens-ai/Pure-Thought/blob/main/PRDs/21-Quantum-LDPC-Codes.md}{PRD
21}. Covered by canonical revision.

\subsection{21. Quantum LDPC instances with independently certified code
parameters}

\textbf{Question.} For a specified CSS construction, can code dimension
and distance bounds be verified independently, and can decoder
performance be measured under one explicit noise model?

\textbf{Why it matters.} Useful quantum computation requires reliable
logical information despite noisy components. Sparse checks offer a
route to distributing error correction efficiently, but code parameters
and decoder behavior must be demonstrated rather than assumed. Exact
small examples are valuable for detecting mistakes that large
simulations can conceal.

\textbf{Scope and model.} Start with hypergraph-product codes from
supplied binary parity-check matrices H\ensuremath{{}_{1}} and
H\ensuremath{{}_{2}}. Include the product of two full-rank
3\ensuremath{\times}7 Hamming parity-check matrices as a small
benchmark. Separate finite instances, asymptotic families, and noisy
syndrome extraction.

Use independent X errors with perfect syndrome measurements for the
first decoder experiment. A family is LDPC only when both row and column
weights stay bounded as n grows.

\subsubsection{Mathematical formulation and claim boundary}

\href{https://github.com/rick-stevens-ai/Pure-Thought/blob/main/PRDs/21-Quantum-LDPC-Codes.md}{Canonical
specification}.

\[H_X=[H_1\otimes I_{n_2}\mid I_{m_1}\otimes H_2^{\mathsf T}],
\qquad H_Z=[I_{n_1}\otimes H_2\mid H_1^{\mathsf T}\otimes I_{m_2}],\]

\[H_XH_Z^{\mathsf T}=0,\quad n=n_1n_2+m_1m_2,\quad
k=n-\operatorname{rank}_{\mathbb F_2}H_X-\operatorname{rank}_{\mathbb F_2}H_Z,\]

\[d_X=\min_{x\in\ker H_Z\setminus\operatorname{row}H_X}|x|,
\qquad d_Z=\min_{z\in\ker H_X\setminus\operatorname{row}H_Z}|z|,
\qquad d=\min(d_X,d_Z).\]

Two full-rank \(3\times7\) Hamming checks give \(n=58\), \(k=16\); the
included independent validator also checks \(d=3\). A logical witness
gives an upper distance bound; exclusion of all smaller logicals gives
the lower bound. Finite code parameters do not establish an asymptotic
LDPC family or decoding threshold.

\textbf{Deliverables.} Export sparse matrices, GF(2) row reductions,
logical bases, explicit upper witnesses and lower-bound proofs. SAT
certificates need a verified encoding of nontrivial homology and weight.
Decoder experiments report physical noise, recovery rule,
logical-failure test, sample counts, confidence intervals, and runtime.

\textbf{Validation.} For the two 3\ensuremath{\times}7 inputs verify
n=58 and k=16; certify the claimed distance independently rather than
assuming it. Inject all small-weight errors on tiny codes. A matching
syndrome is not success unless error plus recovery is a stabilizer.
Compare optimized checks with an independent GF(2) implementation.

\textbf{Milestones.} Baseline: exact small product code and certified
distance bracket. Strong: several instances plus honest finite-size
decoder curves. Research extension: a new family tradeoff or proved
decoder guarantee. Reproducing asymptotically good codes is not itself a
new existence result.

\href{https://github.com/rick-stevens-ai/Pure-Thought/blob/main/reviews/21-Quantum-LDPC-Codes.md}{Detailed
critique}.

\section{Synthesis 22. New Bell Inequalities \& Tsirelson-type Bounds}

\textbf{Mapping:}
\href{https://github.com/rick-stevens-ai/Pure-Thought/blob/main/PRDs/22-Bell-Inequalities-Nonlocality.md}{PRD
22}. Covered by canonical revision.

\subsection{22. Bell inequalities with certified local and quantum value
brackets}

\textbf{Question.} For a fixed rational Bell functional in a small
bipartite scenario, can its exact local value and a certified quantum
lower/upper bracket be produced?

\textbf{Why it matters.} Bell inequalities make the difference between
classical local explanations and quantum correlations precise. They also
underlie ways to certify quantum behavior from observed input-output
relationships. Reliable quantum bounds are essential: a numerical
optimization value without a valid upper certificate or achieving
strategy may misstate what quantum theory permits.

\textbf{Scope and model.} Start with two parties, two binary settings
each, and CHSH. Fix normalization and no-signaling affine coordinates.
The target quantum value is the supremum over finite-dimensional
tensor-product strategies; identify its closure when discussing limits.
Finite NPA levels provide outer bounds through the commuting-operator
framework.

Facet discovery and randomness certification are extensions with
separate definitions. Do not equate a finite SDP solution with a
realizable strategy.

\subsubsection{Mathematical formulation and claim boundary}

\href{https://github.com/rick-stevens-ai/Pure-Thought/blob/main/PRDs/22-Bell-Inequalities-Nonlocality.md}{Canonical
specification}.

\[\beta(P)=\sum_{abxy}c_{abxy}P(ab|xy),\quad
\beta_L=\max_{P\text{ deterministic local}}\beta(P),\quad
L_Q\le\beta_Q\le U_Q.\]

A finite-dimensional state and measurements certify \(L_Q\); a verified
NPA dual or operator SOS certifies \(U_Q\). For CHSH, \(\beta_L=2\) and
an achieving two-qubit strategy plus an upper-bound identity give
\(\beta_Q=2\sqrt2\). NPA's commuting-operator limit must not be
identified with the closure of finite-dimensional tensor-product
strategies. A facet in affine dimension \(D\) needs saturating affine
rank \(D-1\).

\textbf{Deliverables.} Export coefficients, local vertices,
state/measurements, moment-word list, algebraic relations, dual
matrices/SOS, and exact or interval objective bounds. To certify a facet
in affine dimension D, verify the inequality for every vertex and show
saturating vertices have affine dimension D\ensuremath{-}1. Declare
relabeling symmetries when counting classes.

\textbf{Validation.} Check all 16 CHSH deterministic strategies and
\ensuremath{\beta}\_L=2. Reproduce 2\ensuremath{\surd}2 with an explicit
two-qubit strategy and an operator upper-bound identity. Reject an
invalid dual PSD matrix or a wrong commutation relation. Distinguish
eight CHSH variants from one class under relabeling; positivity facets
must be handled separately.

\textbf{Milestones.} Baseline: exact CHSH lower and upper certificates.
Strong: a preselected additional rational functional with a certified
bracket, whether or not it closes. Research extension: a new inequality
with normalized comparison, or device-independent entropy bounds with an
explicit adversary and statistical model.

\href{https://github.com/rick-stevens-ai/Pure-Thought/blob/main/reviews/22-Bell-Inequalities-Nonlocality.md}{Detailed
critique}.

\section{Synthesis 23. Certified Spectral Gaps for Parent Hamiltonians}

\textbf{Mapping:}
\href{https://github.com/rick-stevens-ai/Pure-Thought/blob/main/supplemental/S07.md}{S07}.
Distinct objective preserved in supplement.

\subsection{S07. Finite-size criteria for parent-Hamiltonian spectral
gaps}

\textbf{Question.} For a specified frustration-free one-dimensional
parent-Hamiltonian family, can a validated local spectral calculation
imply a uniform many-body gap?

\textbf{Why it matters.} The gap helps distinguish low-energy behavior
and enters rigorous statements about correlations and stability under
suitable conditions. Proving a uniform gap turns a finite numerical
observation into a statement about a many-body family, which is often
the scientifically relevant object.

\textbf{Scope and model.} Fix exact MPS tensors, on-site dimension,
interaction range, boundary convention, and local positive projectors
h\_i. Verify frustration freeness and ground-space assumptions. Select a
precise Knabe-type or martingale theorem appropriate to that model.
Compute interval enclosures for the required finite-segment nonzero gaps
and verify the theorem's threshold to derive a uniform lower bound for
the stated sizes/boundaries. Degenerate ground spaces require
identifying the first strictly positive eigenvalue rather than assuming
a unique ground state.

\subsubsection{Mathematical formulation and claim boundary}

\href{https://github.com/rick-stevens-ai/Pure-Thought/blob/main/supplemental/S07.md}{Supplemental
specification}.

\[H_N=\sum_i h_i,\qquad h_i\succeq0,\qquad
\Delta_N=\min(\operatorname{spec}H_N\setminus\{0\}),\]

\[\text{target: }\inf_{N\ \mathrm{in\ the\ stated\ family}}\Delta_N\ge\delta>0.\]

Verify frustration freeness and ground-space assumptions. A specified
finite-size gap theorem must supply the implication from validated local
gaps to the displayed uniform bound. Positive gaps on finitely many
chains are insufficient; failure of a sufficient criterion is not a
gaplessness proof.

\textbf{Critique.} The synthesis's jump from MPS to arbitrary PEPS makes
a tractable finite-size criterion look like a general spectral-gap
decision method. Such general algorithms are obstructed by
undecidability results. A positive finite-chain gap alone says nothing
uniform as size grows; one must meet a theorem's hypotheses and
threshold. The revision targets certificate automation for an explicit
restricted family and preserves unresolved outcomes when the chosen
sufficient criterion fails.

\textbf{Validation and milestones.} Baseline: certify a standard parent
model with a valid finite-size theorem. Strong: a uniform parameter
interval with an explicit gap lower bound. Test a deliberately
gapless/threshold-failing control without claiming that criterion
failure proves gaplessness.

Use the shared
\href{https://github.com/rick-stevens-ai/Pure-Thought/blob/main/reviews/EVIDENCE-STANDARD.md}{evidence
standard}. Unsuccessful search is unresolved unless a complete exclusion
proof is supplied. These are proposed baseline specifications, not
completed scientific results.

\href{https://arxiv.org/abs/1502.04573}{Cubitt, Pérez-García and Wolf,
Undecidability of the Spectral Gap} motivates the restriction to a
specified certifiable family.

\section{Synthesis 24. Topological Order \& Anyon Condensation with
Proofs}

\textbf{Mapping:}
\href{https://github.com/rick-stevens-ai/Pure-Thought/blob/main/PRDs/24-Topological-Quantum-Error-Correction.md}{PRD
24},
\href{https://github.com/rick-stevens-ai/Pure-Thought/blob/main/supplemental/S08.md}{S08}.
Related PRD; supplement preserves distinct objective.

\subsection{24. Topological code verification and noise-specific decoder
evaluation}

\textbf{Question.} Can a toric-code family's algebraic parameters be
proved, and can a fixed decoder's logical failure probability be
evaluated under a precisely specified error model?

\textbf{Why it matters.} Geometric codes connect an intuitive picture of
local errors to precise algebraic guarantees. They are important models
for protecting quantum information using structured measurements.
Understanding logical failure, rather than only syndrome cancellation,
is essential for assessing whether increasing the code size actually
improves reliability.

\textbf{Scope and model.} Use an L\ensuremath{\times}L square torus with
qubits on edges and L\ensuremath{\ge}3. Start with independent Z errors
of probability p and perfect syndrome extraction. Specify the decoder,
tie-breaking rule, and whether recovery weights are uniform or
likelihood based.

Keep exact code properties separate from statistical decoder
performance. Planar boundaries, color codes, measurement noise, and
circuit-level noise are distinct extensions.

\subsubsection{Mathematical formulation and claim boundary}

\href{https://github.com/rick-stevens-ai/Pure-Thought/blob/main/PRDs/24-Topological-Quantum-Error-Correction.md}{Canonical
specification}.

\[C_2\xrightarrow{\partial_2}C_1\xrightarrow{\partial_1}C_0,
\quad\partial_1\partial_2=0,\quad H_X=\partial_1,\quad H_Z=\partial_2^{\mathsf T},\]

\[n=2L^2,\quad k=2,\quad d=L\quad(L\ge3),\qquad
e+r\in\operatorname{im}\partial_2\quad\text{for successful }Z\text{-error recovery}.\]

A vanishing residual syndrome only gives \(e+r\in\ker\partial_1\) and
can hide a nontrivial logical cycle. Statistical decoder curves require
the fixed noise model and decoder; finite-size crossings are threshold
estimates, not limit theorems.

\textbf{Deliverables.} Export lattice incidence data, ranks and logical
bases, distance argument, decoder specification, and failure counts with
confidence intervals. Exhaustive small-size evaluation can give an exact
failure polynomial in p.~Monte Carlo output must retain trials and
failure counts, including zero-failure upper intervals.

\textbf{Validation.} Check n,k,d on several small tori. Verify logical
anticommutation via primal/dual intersection. Enumerate all low-weight
errors on tiny instances and test the stabilizer-membership success
criterion. Validate a matching solution's optimality separately from
logical decoding success.

\textbf{Milestones.} Baseline: exact construction and distance proof
with a tested decoder. Strong: reproducible finite-size failure curves
and uncertainty-aware threshold estimates. Research extension: a
rigorous threshold bound or noisy-measurement/circuit model, with no
transfer of numerical thresholds across noise conventions.

\href{https://github.com/rick-stevens-ai/Pure-Thought/blob/main/reviews/24-Topological-Quantum-Error-Correction.md}{Detailed
critique}.

\subsection{S08. Gapped boundaries from exact anyon data}

\textbf{Question.} For one fully specified two-dimensional topological
order, can candidate condensates and their gapped boundaries be verified
from exact categorical data?

\textbf{Why it matters.} Boundaries determine how a topological phase
interfaces with its surroundings and which excitations can disappear or
remain distinguishable there. Understanding these options supports the
theory of topological phases and the organization of quantum information
in such systems.

\textbf{Scope and model.} Start with the toric-code modular category,
explicit simple objects, fusion, twists, F and R symbols and
conventions. Enumerate candidate condensable algebra objects within a
declared multiplicity bound. Check multiplication/unit, associativity,
commutativity and separability conditions, and the Lagrangian dimension
condition for a boundary to vacuum. For the abelian toric-code benchmark
recover electric and magnetic condensates. If constructing a lattice
boundary, supply its Hamiltonian and check commuting-projector and gap
claims separately.

\subsubsection{Mathematical formulation and claim boundary}

\href{https://github.com/rick-stevens-ai/Pure-Thought/blob/main/supplemental/S08.md}{Supplemental
specification}.

\[A=\bigoplus_a n_a a,\qquad n_a\in\mathbb Z_{\ge0},\qquad
(\dim A)^2=\operatorname{Dim}\mathcal C.\]

For a vacuum boundary this Lagrangian dimension condition supplements,
rather than replaces, the unit, multiplication, associativity,
commutativity and separability checks using exact categorical coherence
data. In the toric-code category the reference condensates are
\(1\oplus e\) and \(1\oplus m\). A microscopic boundary Hamiltonian
needs its own gap proof.

\textbf{Critique.} Fusion rules and modular S,T tables alone do not in
general specify or prove all categorical coherence data. The synthesis
also mixes 2D anyons with 3D topological orders, which require different
structures. A listed condensate is not automatically a microscopic
boundary realization. The rewrite retains an exact categorical task,
bounds the enumeration and separates algebraic classification from
lattice realization and anomaly/inflow interpretation.

\textbf{Validation and milestones.} Baseline: verify the toric-code
coherence data and two standard vacuum boundaries. Strong: a second
fully specified category or a completeness proof within a bounded class.
Invalid F/R symbols and a nonbosonic condensate are rejection controls.

Use the shared
\href{https://github.com/rick-stevens-ai/Pure-Thought/blob/main/reviews/EVIDENCE-STANDARD.md}{evidence
standard}. Unsuccessful search is unresolved unless a complete exclusion
proof is supplied. These are proposed baseline specifications, not
completed scientific results.

\href{https://arxiv.org/abs/2504.19512}{Li, Yang and Dong, Gapped
Boundaries of Kitaev's Quantum Double Models} provides a primary example
of linking Lagrangian algebras to explicit lattice boundary terms.

\section{Synthesis 25. Sign-Problem (Un)avoidability}

\textbf{Mapping:}
\href{https://github.com/rick-stevens-ai/Pure-Thought/blob/main/supplemental/S09.md}{S09}.
Distinct objective preserved in supplement.

\subsection{S09. Stoquasticity under a bounded class of basis changes}

\textbf{Question.} For a finite rational Pauli Hamiltonian, is there an
allowed product Clifford transformation making all off-diagonal entries
real and nonpositive?

\textbf{Why it matters.} Knowing when signs can be removed helps
separate difficulties caused by representation from those that survive a
proposed remedy. A bounded search also creates small, exact benchmarks
for algorithms intended to find useful basis changes.

\textbf{Scope and model.} Fix n qubits, rational coefficients, and
allowed U=\ensuremath{\otimes}\_iU\_i with U\_i from the finite
single-qubit Clifford group modulo phase. Define stoquasticity in the
computational basis for the full transformed H, not automatically term
by term. Exhaustively enumerate or use an exactly equivalent constraint
encoding. Export U and an exact matrix/sign check if successful;
otherwise export complete enumeration or a checked UNSAT proof.
Continuous local unitaries require a semialgebraic formulation and are a
distinct extension.

\subsubsection{Mathematical formulation and claim boundary}

\href{https://github.com/rick-stevens-ai/Pure-Thought/blob/main/supplemental/S09.md}{Supplemental
specification}.

\[U=\bigotimes_{i=1}^nU_i,\quad U_i\in\mathrm{Clifford}_1/\mathrm{phase},
\qquad H'=UHU^\dagger,\]

\[\operatorname{Im}H'_{xy}=0,\qquad\operatorname{Re}H'_{xy}\le0\quad(x\ne y).\]

Apply these conditions to the full transformed matrix, allowing
cancellations between terms. Exhaustive coverage or an equivalent
checked UNSAT proof excludes only this finite transformation class.
Stoquasticity alone does not guarantee efficient mixing or sampling.

\textbf{Critique.} The original omits which local basis changes are
allowed, so ``unavoidable sign problem'' has no definite quantifier. A
SAT search over finitely many transformations cannot exclude all
continuous local rotations. Requiring every term to be stoquastic can
also be stronger than requiring their sum to be stoquastic because
entries can cancel. Finally, stoquasticity does not guarantee efficient
sampling or mixing. The rewrite makes a finite, falsifiable basis-cure
problem and limits every no-go conclusion to its transformation class.

\textbf{Validation and milestones.} Baseline: exhaustive small-n
decisions with explicit transformations or coverage. Strong: verified
SAT encoding checked against brute force. Extend to a continuous
rotation class only with exact equations, inequalities, and global proof
machinery.

Use the shared
\href{https://github.com/rick-stevens-ai/Pure-Thought/blob/main/reviews/EVIDENCE-STANDARD.md}{evidence
standard}. Unsuccessful search is unresolved unless a complete exclusion
proof is supplied. These are proposed baseline specifications, not
completed scientific results.

\href{https://pmc.ncbi.nlm.nih.gov/articles/PMC6450938/}{Marvian, Lidar
and Hen, On the computational complexity of curing non-stoquastic
Hamiltonians} supports the need to specify the allowed basis class.

\section{Synthesis 26. KAM/Nekhoroshev Stability Domains for Planetary
Systems}

\textbf{Mapping:}
\href{https://github.com/rick-stevens-ai/Pure-Thought/blob/main/PRDs/26-KAM-Theory-Planetary-Stability.md}{PRD
26},
\href{https://github.com/rick-stevens-ai/Pure-Thought/blob/main/PRDs/28-Nekhoroshev-Stability-Theory.md}{PRD
28}. Covered by canonical revision.

\subsection{26. A posteriori certification of an invariant Hamiltonian
torus}

\textbf{Question.} For an explicitly specified analytic near-integrable
Hamiltonian and an approximate torus, can every hypothesis of a
quantitative KAM theorem be verified?

\textbf{Why it matters.} Long-term stability is difficult to infer from
a finite trajectory simulation. A theorem can instead prove that a
particular organized motion exists despite perturbations. This offers a
foundation for studying stability mechanisms in celestial and other
Hamiltonian systems, provided the model and theorem assumptions are
appropriate.

\textbf{Scope and model.} Begin with
H(I,\ensuremath{\theta})=\ensuremath{\omega}\ensuremath{\cdot}I+|I|\ensuremath{{}^{2}}/2+\ensuremath{\epsilon}[cos
\ensuremath{\theta}\ensuremath{{}_{1}}+cos(\ensuremath{\theta}\ensuremath{{}_{1}}\ensuremath{-}\ensuremath{\theta}\ensuremath{{}_{2}})]
in two degrees of freedom,
\ensuremath{\omega}=(1,(1+\ensuremath{\surd}5)/2), on a stated complex
neighborhood of a real action box. Choose a rational
\ensuremath{\epsilon} only after validating the unperturbed case. Seek
an embedding with rotation vector \ensuremath{\omega}.

This is a mathematical model, not a solar-system stability assertion.
Planetary extensions require an explicit canonical reduction, handling
Kepler degeneracy, collision exclusions, and input uncertainty.

\subsubsection{Mathematical formulation and claim boundary}

\href{https://github.com/rick-stevens-ai/Pure-Thought/blob/main/PRDs/26-KAM-Theory-Planetary-Stability.md}{Canonical
specification}.

\[H(I,\theta)=\omega\cdot I+\frac{|I|^2}{2}
+\epsilon[\cos\theta_1+\cos(\theta_1-\theta_2)],
\qquad\omega=(1,(1+\sqrt5)/2),\]

\[E(\theta)=X_H(K(\theta))-DK(\theta)\omega,
\qquad |k\cdot\omega|\ge\frac\gamma{|k|_1^\tau}
\quad(k\in\mathbb Z^2\setminus\{0\}).\]

Prove the all-integer Diophantine bound and every analytic-width, twist,
tail and smallness condition of a specified quantitative theorem. Its
conclusion is a nearby invariant torus, not stability of arbitrary
planetary initial conditions. Failure of a sufficient criterion is not
instability.

\textbf{Deliverables.} Export H, \ensuremath{\epsilon}, domain, exact
frequency representation, Fourier coefficients, all theorem constants,
interval inequalities, and a checker. State the conclusion: existence of
a nearby quasiperiodic invariant torus with a quantified embedding
error. A failed sufficient condition yields ``not certified,'' not
``unstable.''

\textbf{Validation.} Verify E=0 for
K\ensuremath{{}_{0}}(\ensuremath{\theta})=(I=0,\ensuremath{\theta}) at
\ensuremath{\epsilon}=0. Derive an all-integer Diophantine bound for the
golden-ratio frequency. Check Fourier tails and loss of analytic width.
Test a resonant frequency where the stated theorem must refuse
certification.

\textbf{Milestones.} Baseline: unperturbed certificate and complete
arithmetic/theorem setup. Strong: a nonzero-\ensuremath{\epsilon} torus
certificate. Research extension: a larger perturbation range or a
reduced planetary model. One invariant torus does not certify a
neighborhood of arbitrary observed initial conditions or a phase-space
measure bound.

\href{https://github.com/rick-stevens-ai/Pure-Thought/blob/main/reviews/26-KAM-Theory-Planetary-Stability.md}{Detailed
critique}.

\subsection{28. Explicit finite-time action confinement from a normal
form}

\textbf{Question.} For a specified analytic Hamiltonian and action
domain, what rigorously computable action excursion and confinement time
follow from a certified normal form or quantitative Nekhoroshev theorem?

\textbf{Why it matters.} Finite-time guarantees answer a question that
simulations alone cannot settle beyond their run length. They also
distinguish a valid stability argument from an impressive-looking
exponential formula with missing constants or units. Understanding slow
drift is relevant to the mathematical study of planetary and other
nearly integrable systems.

\textbf{Scope and model.} Use
H=h(I)+\ensuremath{\epsilon}f(I,\ensuremath{\theta}),
h=\ensuremath{\omega}\ensuremath{\cdot}I+|I|\ensuremath{{}^{2}}/2, in
two or three degrees of freedom on a compact action box with an explicit
complex extension. Supply \ensuremath{\omega}, a finite Fourier
polynomial f, and rational \ensuremath{\epsilon}. Bound the gradient
away from zero where the selected theorem requires it.

The initial target is a valid dimensional time bound, even if modest.
Planetary and post-Newtonian models are extensions requiring their own
reductions and hypotheses.

\subsubsection{Mathematical formulation and claim boundary}

\href{https://github.com/rick-stevens-ai/Pure-Thought/blob/main/PRDs/28-Nekhoroshev-Stability-Theory.md}{Canonical
specification}.

\[H\circ\Phi=Z(J)+R(J,\theta),\qquad
|I-J|\le\eta,\qquad\|\partial_\theta R\|\le r,\]

\[|I(t)-I(0)|\le2\eta+r|t|.\]

This estimate assumes the angle-independent normal form shown and a
proved domain-containment interval. Include resonant drift if \(Z\)
depends on angles. A quantitative Nekhoroshev theorem instead requires
all its constants, hypotheses and threshold in
\(T=C\exp(c\epsilon^{-a})\), with units supplied by \(C\). The number
\(\exp(10^{0.3})\approx7.35413\) has no unit of time.

\textbf{Deliverables.} Export normal-form order, transformations,
inverse/domain bounds, Fourier tails, r, \ensuremath{\eta}, and a
domain-containment proof for the claimed interval. A Nekhoroshev result
must include all constants and its smallness threshold in T=C exp(c
\ensuremath{\epsilon}\textasciicircum{}\ensuremath{-}a) with stated
units; do not choose a from a desired time target.

\textbf{Validation.} Verify integrable \ensuremath{\epsilon}=0 behavior
and a low-order symbolic normal form. Check symplecticity with validated
remainder control. Confirm the action bound stays inside the certified
domain. Independently recompute
exp(10\textasciicircum{}0.3)\ensuremath{\approx}7.35: it is a
dimensionless factor, not 10\ensuremath{{}^{13}} years.

\textbf{Milestones.} Baseline: one explicit finite-time confinement
certificate. Strong: optimized truncation with a stronger bound or full
quantitative quasi-convex theorem application. Research extension:
exponential scaling across a parameter family; do not infer actual
solar-system stability from \ensuremath{\epsilon} alone.

\href{https://github.com/rick-stevens-ai/Pure-Thought/blob/main/reviews/28-Nekhoroshev-Stability-Theory.md}{Detailed
critique}.

\section{Synthesis 27. Periodic Orbits \& Invariant Manifolds in N-Body
Problem}

\textbf{Mapping:}
\href{https://github.com/rick-stevens-ai/Pure-Thought/blob/main/supplemental/S10.md}{S10}.
Distinct objective preserved in supplement.

\subsection{S10. Certified periodic orbits of a specified three-body
model}

\textbf{Question.} Can a numerically proposed periodic orbit be enclosed
by a validated shooting proof and classified by its reduced Floquet
spectrum?

\textbf{Why it matters.} Periodic orbits are organizing reference
motions in gravitational dynamics. Certifying their existence gives
subsequent studies of nearby motion a dependable starting point and
separates exact solutions from visual patterns in a simulation.

\textbf{Scope and model.} Choose one nondimensional Newtonian three-body
model with positive rational masses, or a planar circular restricted
model with fixed mass ratio. Declare collision exclusions and conserved
quantities. Define a Poincaré section or phase condition to remove time
translation, and reduce continuous symmetries. Solve the resulting
shooting residual using a validated flow integrator and interval
Newton/Krawczyk bounds for initial conditions and period. Integrate
variational equations with enclosures for the monodromy matrix; remove
the expected neutral directions before interpreting multipliers.

\subsubsection{Mathematical formulation and claim boundary}

\href{https://github.com/rick-stevens-ai/Pure-Thought/blob/main/supplemental/S10.md}{Supplemental
specification}.

\[F(z_0,T)=\varphi_T(z_0)-z_0=0,\qquad
M=D\varphi_T(z_0).\]

The shooting residual must first be augmented and reduced using a phase
condition, conserved quantities and continuous symmetries so the
validated root problem has the correct dimension. Interval flow and root
bounds establish existence; the reduced Floquet multipliers classify
linear behavior. Neither a near return nor stable-looking multipliers
prove nonlinear stability or a manifold connection.

\textbf{Critique.} The original bundles orbit existence, complete family
discovery, linear stability, invariant manifolds and heteroclinic
connections into one certificate. A small shooting residual is not
existence, and stable-looking Floquet multipliers are not nonlinear
stability. Near-return trajectories do not establish a connection
between invariant manifolds. The revised first target is a single orbit
with a phase-fixed existence proof; continuation, manifold enclosure and
connection proofs are explicit later stages. This also preserves the
figure-eight idea incorrectly folded into PRD 27's central
configurations.

\textbf{Validation and milestones.} Baseline: certify one supplied orbit
and its period. Strong: validated continuation over a parameter interval
and a reduced Floquet classification. A manifold connection requires
additional covering/intersection proof, not a plot of approaching
trajectories.

Use the shared
\href{https://github.com/rick-stevens-ai/Pure-Thought/blob/main/reviews/EVIDENCE-STANDARD.md}{evidence
standard}. Unsuccessful search is unresolved unless a complete exclusion
proof is supplied. These are proposed baseline specifications, not
completed scientific results.

\section{Synthesis 28. Central Configurations Classification}

\textbf{Mapping:}
\href{https://github.com/rick-stevens-ai/Pure-Thought/blob/main/PRDs/27-N-Body-Central-Configurations.md}{PRD
27}. Covered by canonical revision.

\subsection{27. Central configurations in a bounded, symmetry-reduced
N-body problem}

\textbf{Question.} For fixed positive masses and a declared symmetry
quotient, can all collision-free central configurations in the chosen
family be certified and their relative-equilibrium linearization
analyzed?

\textbf{Why it matters.} Central configurations are organizing
structures in gravitational dynamics and provide exact reference cases
amid otherwise complicated motion. A complete small catalog can test
numerical methods and reveal how mass choices and symmetry affect
possible motions. The problem also links gravity to exact algebraic root
finding.

\textbf{Scope and model.} Start with three labeled positive rational
masses in the plane. Fix G=1, center of mass zero, moment of inertia
I=\ensuremath{\Sigma}m\_i|r\_i|\ensuremath{{}^{2}}=1, and quotient by
O(2), so mirror images count together. Permute only equal masses when
explicitly requested. Extend to a specified four-body symmetry family
after validating the complete three-body case.

Central configurations generate homographic motions; general periodic
choreographies are a separate problem.

\subsubsection{Mathematical formulation and claim boundary}

\href{https://github.com/rick-stevens-ai/Pure-Thought/blob/main/PRDs/27-N-Body-Central-Configurations.md}{Canonical
specification}.

\[U=\sum_{i<j}\frac{m_im_j}{|r_i-r_j|},\quad
\sum_i m_i r_i=0,\quad I=\sum_i m_i|r_i|^2=1,\]

\[\nabla_{r_i}U=-\lambda m_i r_i,\qquad\lambda=U/I>0,
\qquad d_{ij}^2=|r_i-r_j|^2,\quad d_{ij}>0.\]

Preserve collision exclusions when clearing denominators. For three
labeled positive masses modulo \(O(2)\), recover three collinear
orderings and one equilateral shape. Complete coverage requires more
than isolated root boxes. Relative-equilibrium stability uses the
reduced rotating-frame system including Coriolis terms; a figure-eight
choreography is not a central configuration.

\textbf{Deliverables.} Return normalized configurations, distance
enclosures, residual identities, existence/uniqueness proofs,
equivalence representatives, and completeness scope. For planar relative
equilibria analyze the reduced rotating-frame first-order system
including Coriolis terms. Distinguish spectral stability from nonlinear
stability.

\textbf{Validation.} For three labeled masses under the stated O(2)
quotient, recover three collinear orderings (one per middle body) and
one equilateral shape. Check \ensuremath{\lambda}=U at I=1. Test
collision artifacts introduced by denominator clearing. For stability,
reproduce a known three-body criterion using identical mass and
reduction conventions.

\textbf{Milestones.} Baseline: complete certified three-body catalog.
Strong: complete enumeration in one four-body symmetry family or mass
instance, with all charts covered. Research extension: a five-body
restricted family; no claim to solve the general finiteness problem
without an appropriately global proof.

\href{https://github.com/rick-stevens-ai/Pure-Thought/blob/main/reviews/27-N-Body-Central-Configurations.md}{Detailed
critique}.

\section{Synthesis 29. Minimal Autocatalytic Cores \& Universality}

\textbf{Mapping:}
\href{https://github.com/rick-stevens-ai/Pure-Thought/blob/main/PRDs/29-Chemical-Reaction-Networks-Origins.md}{PRD
29},
\href{https://github.com/rick-stevens-ai/Pure-Thought/blob/main/supplemental/S11.md}{S11}.
Distinct objective preserved in supplement.

\subsection{29. Autocatalytic sets: maximality, irreducibility and
dynamical viability}

\textbf{Question.} For a finite catalytic reaction system, can its
maximal RAF and smallest RAF subsets be certified, while keeping
structural autocatalysis distinct from sustained growth?

\textbf{Why it matters.} Collectively supported chemistry is relevant to
thinking about how organized reaction systems can arise and persist.
Exact minimality results help distinguish essential reactions from
dispensable ones in a model. Just as importantly, separating structural
closure from energy supply and kinetics prevents a network diagram from
being mistaken for proof of a viable replicating system.

\textbf{Scope and model.} Input species S, reactions R with integer
stoichiometry, food set F, and catalysis relation
C\ensuremath{\subset}S\ensuremath{\times}R. Begin with at most 12
reactions for exhaustive cross-checks. Catalysis is supplied model data,
not inferred from a stoichiometric cycle.

Define three separate targets: unique maximal RAF, an inclusion-minimal
(irreducible) RAF, and a minimum-cardinality RAF. Thermodynamic and
kinetic viability require extra inputs and are separate extensions.

\subsubsection{Mathematical formulation and claim boundary}

\href{https://github.com/rick-stevens-ai/Pure-Thought/blob/main/PRDs/29-Chemical-Reaction-Networks-Origins.md}{Canonical
specification}.

\[R'\subseteq R,\quad R'\ne\varnothing,\quad
\operatorname{reactants}(r)\subseteq\operatorname{cl}_{R'}(F),\]

\[\forall r\in R'\ \exists s\in\operatorname{cl}_{R'}(F):\ (s,r)\in C.\]

Compute closure by reactant availability without imposing catalysis in
the closure step. These conditions define a RAF. Inclusion-maximal,
inclusion-minimal (irreducible), and minimum-cardinality RAFs are
separate targets. Catalysis-respecting startup and sustained growth
require additional conditions.

\textbf{Deliverables.} Export closure rounds, catalyst witnesses, maxRAF
elimination trace, and minimality/cardinality proof with explicit scope.
For thermodynamic extension specify chemical potentials, chemostats, and
driving. For kinetic extension provide rates, dilution, initial
conditions and a positive sustained-state/growth criterion; RAF
existence alone is insufficient.

\textbf{Validation.} Cross-check all subsets for tiny networks. Include
food-catalyzed, self-catalyzed, catalyst-free, and mutually dependent
examples. Verify a maximum RAF can contain several irreducible RAFs of
different sizes. Deletion tests for irreducibility must use a complete
RAF detector on each remaining network, not merely test whether the
entire remainder is itself RAF.

\textbf{Milestones.} Baseline: checked maxRAF/RAF witnesses. Strong:
complete minimum-cardinality results for a bounded network family.
Research extension: characterize which RAFs additionally support a
declared driven mass-action regime. Biological self-replication and
evolution require their own operational definitions.

\href{https://github.com/rick-stevens-ai/Pure-Thought/blob/main/reviews/29-Chemical-Reaction-Networks-Origins.md}{Detailed
critique}.

\subsection{S11. Minimal reaction implementations of a specified
computational behavior}

\textbf{Question.} Within a bounded reaction grammar, what is the
smallest network implementing a precisely defined input/output behavior
with a proved error bound?

\textbf{Why it matters.} Chemical computation explores how local
reactions can perform information processing, including ideas relevant
to molecular sensing and control. Precise semantics make proposed
constructions comparable and prevent hidden resources or flexible
readout choices from creating misleading claims of simplicity.

\textbf{Scope and model.} Choose deterministic mass-action or stochastic
reaction semantics, species/reaction bounds, stoichiometric arity, input
encoding, output decoding and cost (species count first, then reaction
count, for example). Begin with a finite Boolean input/output function
and an explicit settling-time/error requirement. Enumerate networks
modulo species relabeling; prove correctness over the complete allowed
input set and exclude every lower-cost network within the grammar. For a
universal-computation extension, supply a uniform compiler and a
simulation theorem covering arbitrary program lengths, resource scaling
and precision.

\subsubsection{Mathematical formulation and claim boundary}

\href{https://github.com/rick-stevens-ai/Pure-Thought/blob/main/supplemental/S11.md}{Supplemental
specification}.

\[\min_{\mathcal R\in\mathfrak G}\operatorname{cost}(\mathcal R)
\quad\text{subject to}\quad
\forall x\in\mathcal X:\ \operatorname{error}(\mathcal R,x;T)\le\varepsilon.\]

Declare the finite reaction grammar \(\mathfrak G\), semantics,
encoding, decoding, time requirement, input set and cost ordering.
Correctness of one network gives an upper cost bound; excluding every
lower-cost network gives minimality within that grammar. A finite truth
table does not prove universal computation or self-replication.

\textbf{Critique.} The synthesis joins minimal autocatalysis,
self-replication and universal computation, but none implies the others.
``Smallest chemistry'' depends on kinetics, encoding, concentration
precision and which resources count as free. A fixed finite truth table
does not prove Turing universality, while an analog model with unbounded
precision may hide computation in its inputs. The rewrite separates a
bounded minimal implementation result from a universal simulation
theorem and from the RAF structure in PRD 29.

\textbf{Validation and milestones.} Baseline: one certified reaction
implementation with explicit semantics. Strong: minimality within a
finite grammar and input domain. Extension: a compiler/simulation
theorem; avoid calling the baseline universal or biologically
self-replicating.

Use the shared
\href{https://github.com/rick-stevens-ai/Pure-Thought/blob/main/reviews/EVIDENCE-STANDARD.md}{evidence
standard}. Unsuccessful search is unresolved unless a complete exclusion
proof is supplied. These are proposed baseline specifications, not
completed scientific results.

\section{\texorpdfstring{Synthesis 30. Genotype\ensuremath{\to}Phenotype
Channel Capacity Bounds}{Synthesis 30.  Channel Capacity Bounds}}

\textbf{Mapping:}
\href{https://github.com/rick-stevens-ai/Pure-Thought/blob/main/PRDs/30-Genotype-Phenotype-Mapping.md}{PRD
30},
\href{https://github.com/rick-stevens-ai/Pure-Thought/blob/main/supplemental/S12.md}{S12}.
Related PRD; supplement preserves distinct objective.

\subsection{30. An exact finite genotype--phenotype map and its mutation
graph}

\textbf{Question.} For a fully specified combinatorial RNA folding rule
on short sequences, what are the exact neutral components, mutation
robustness, and accessible phenotype counts?

\textbf{Why it matters.} The structure of a genotype--phenotype map
helps explain how variation can accumulate without immediately changing
a selected outcome, and how new outcomes become accessible. Exact small
maps provide benchmarks for sampling methods and expose how conclusions
depend on folding rules and tie choices.

\textbf{Scope and model.} Use alphabet \{A,C,G,U\}, sequence length
L=4,\ldots,8 initially, allowed pairs AU/UA, GC/CG and GU/UG, minimum
hairpin length three unpaired bases, no pseudoknots, and one pair per
base. Maximize pair count and break ties by lexicographically smallest
dot-bracket string under a declared character order.

This defines a mathematical GP map, not a thermodynamic RNA predictor. A
nearest-neighbor free-energy extension must name its parameter set;
empirical parameters are then acknowledged model inputs.

\subsubsection{Mathematical formulation and claim boundary}

\href{https://github.com/rick-stevens-ai/Pure-Thought/blob/main/PRDs/30-Genotype-Phenotype-Mapping.md}{Canonical
specification}.

\[\Phi:\{A,C,G,U\}^L\longrightarrow\mathcal S_L,\qquad
r(s)=\frac{|\{t:d_H(s,t)=1,\ \Phi(t)=\Phi(s)\}|}{3L},\]

\[\sum_{\varphi}|\Phi^{-1}(\varphi)|=4^L,\qquad
\deg_{H(L,4)}s=3L.\]

Specify pair-count maximization, allowed pairs, hairpin length,
pseudoknot exclusion and deterministic tie-breaking. Enumerate every
connected component of each phenotype preimage. This combinatorial map
is not a thermodynamic predictor. At \(L=20\),
\(4^L=1{,}099{,}511{,}627{,}776\); sampling does not establish a
complete catalog.

\textbf{Deliverables.} Export folding rules, tie-breaking,
sequence-to-structure table, neutral components, exact counts and
rational robustness. A search restricted to a component must not report
the full phenotype frequency. Fitness is a separately supplied function;
evolutionary dynamics require population size, mutation kernel,
selection rule and initial state.

\textbf{Validation.} Cross-check dynamic programming with exhaustive
legal-structure enumeration on short sequences, especially ties and
prohibited pairs. Verify
\ensuremath{\Sigma}\_\ensuremath{\phi}|\ensuremath{\Phi}\ensuremath{{}^{-1}}(\ensuremath{\phi})|=4\textasciicircum{}L
and every Hamming vertex degree is 3L. Count neutral edges independently
and recover average robustness from twice the edge count. Treat any
larger-L sample as statistical evidence.

\textbf{Milestones.} Baseline: complete exact maps through L=8. Strong:
increase L within measured resources and characterize component-level
robustness/evolvability. Research extension: a proved asymptotic
property or a specified noisy GP-channel capacity problem. Do not
hard-code a universal percolation threshold or empirical target
correlation.

\href{https://github.com/rick-stevens-ai/Pure-Thought/blob/main/reviews/30-Genotype-Phenotype-Mapping.md}{Detailed
critique}.

\subsection{S12. Capacity of a specified finite noisy
genotype--phenotype channel}

\textbf{Question.} For a finite conditional phenotype law P(y\textbar g)
and optional genotype cost constraint, what certified lower and upper
bounds can be placed on Shannon capacity?

\textbf{Why it matters.} An explicit channel separates the number of
possible outputs from their reliability and distinguishability. Capacity
bounds give a rigorous way to compare modeled noise mechanisms or
resource constraints, and to check whether a numerical optimization has
found the best input distribution.

\textbf{Scope and model.} Supply finite genotype and phenotype
alphabets, a stochastic matrix W(y\textbar g), and optional costs c(g)
with E[c]\ensuremath{\le}C. Define capacity max\_p I\_p(G;Y), with
logarithms base two. A chosen input distribution gives a lower bound.
For the unconstrained channel, any output distribution q with adequate
support gives upper bound max\_g D(W(.\textbar g)\textbar\textbar q);
constrained dual bounds must include the cost multiplier. Use interval
logarithms and validated probabilities. Connect to PRD 30 by
constructing W from a fully declared folding/noise rule.

\subsubsection{Mathematical formulation and claim boundary}

\href{https://github.com/rick-stevens-ai/Pure-Thought/blob/main/supplemental/S12.md}{Supplemental
specification}.

\[C=\max_p I_p(G;Y),\qquad
I_p(G;Y)\le C\le\max_g D(W(\cdot|g)\|q).\]

For a cost limit \(c_*\), any \(\lambda\ge0\) and adequately supported
\(q\) give

\[C(c_*)\le\lambda c_*+
\max_g\{D(W(\cdot|g)\|q)-\lambda c(g)\}.\]

Use base-two logarithms and verified stochastic rows. A feasible input
distribution supplies a lower bound, and interval logarithms certify
numerical values. For a deterministic unrestricted finite map, capacity
is \(\log_2|\operatorname{im}\Phi|\); this differs from neutral-network
connectivity.

\textbf{Critique.} Mutual information is a property of a joint
distribution, not a map alone. The original omits the input
distribution, noise, phenotype alphabet and optimization constraints.
For a deterministic finite map with unrestricted input distribution,
capacity is simply log\ensuremath{{}_{2}} of the number of reachable
phenotypes, so a purported hard capacity problem may collapse. Fano/Le
Cam bounds answer specific inference questions and are not automatic
capacity certificates. The rewrite specifies a nontrivial noisy channel,
separates achievability from converse bounds, and keeps biophysical
interpretation conditional on the chosen channel.

\textbf{Validation and milestones.} Baseline: binary symmetric-channel
capacity and deterministic-map capacity reproduced. Strong: a certified
bracket for a finite GP channel and a cost-constrained extension with a
valid dual. Check rows sum to one, handle zero probabilities explicitly,
and report unresolved numerical gaps.

Use the shared
\href{https://github.com/rick-stevens-ai/Pure-Thought/blob/main/reviews/EVIDENCE-STANDARD.md}{evidence
standard}. Unsuccessful search is unresolved unless a complete exclusion
proof is supplied. These are proposed baseline specifications, not
completed scientific results.

\section{Evidence and execution standard}

This standard applies to the revised specifications. It is a proposed
contract for future research runs; this editorial revision does not
supply the scientific certificates themselves.

\subsection{What ``pure thought'' means here}

Use explicitly stated mathematical models and deductions without fitting
conclusions to experimental outcomes. Synthetic examples, exact model
inputs, analytic references, numerical search, and established software
are allowed. Physical constants, electronic-structure approximations,
constitutive laws, and empirical parameter sets must be identified when
used. A model-derived result is conditional on that model.

The original ban on legacy software conflicts with its proposed reliance
on NumPy, SciPy, CVXPY, nauty, and other libraries. Reimplementing every
optimizer and algebra system is not necessary for scientific
independence. Prefer a small independently checkable certificate and a
documented trusted base. A fresh implementation can be a project
objective, but must not be confused with a mathematical proof or novelty
claim.

\subsection{Four evidence levels}

{\def\LTcaptype{none} % do not increment counter
\begin{longtable}[]{@{}
  >{\raggedright\arraybackslash}p{(\linewidth - 4\tabcolsep) * \real{0.3333}}
  >{\raggedright\arraybackslash}p{(\linewidth - 4\tabcolsep) * \real{0.3333}}
  >{\raggedright\arraybackslash}p{(\linewidth - 4\tabcolsep) * \real{0.3333}}@{}}
\toprule\noalign{}
\begin{minipage}[b]{\linewidth}\raggedright
Label
\end{minipage} & \begin{minipage}[b]{\linewidth}\raggedright
What it establishes
\end{minipage} & \begin{minipage}[b]{\linewidth}\raggedright
What it does not establish
\end{minipage} \\
\midrule\noalign{}
\endhead
\bottomrule\noalign{}
\endlastfoot
Numerical evidence & A reproducible computation at stated precision,
sample size, and discretization & Exact truth, continuum coverage, or a
universal theorem \\
Finite-model certificate & An exact/validated statement about the
specified finite model or finite relaxation & An omitted tail,
thermodynamic limit, or physical realization \\
Conditional mathematical theorem & A proof over the stated domain with
every hypothesis and error controlled & Truth outside those assumptions
or a demonstrated physical system \\
Formal proof & A checked theorem in a named proof assistant, with
imported axioms/dependencies listed & Correct physical modeling merely
by virtue of formalization \\
\end{longtable}
}

Formalization is an optional extra deliverable unless the specific
project makes it mandatory. High precision, JSON output, solver status
and small residuals do not by themselves confer any of these proof
levels.

\subsection{Claim directions}

For minimization over a larger feasible set, the optimum is a lower
bound on the original minimum. A physical feasible point supplies an
upper bound. For maximization the directions reverse. Track feasibility
of the primal/dual witnesses separately from their objective values.

Infeasibility of an inner ansatz excludes only that ansatz.
Infeasibility of an outer relaxation can exclude the original problem if
the map into the relaxation is proved and all approximations preserve
necessity. Passing necessary tests means ``not excluded''; existence
needs a constructive sufficient argument.

A smallest observed example is not a minimality theorem. A complete
catalog must define the universe, equivalence relation, boundaries, and
proof that pruning misses nothing. A finite-size numerical trend does
not prove an asymptotic limit.

\subsection{Minimum run manifest}

Before calculation, record the problem ID and revision hash, exact
inputs and conventions, quantified domain, equivalence relation, target
observable, accepted evidence level, tolerances, and resource limits.
Record dependencies and precision/rounding mode. Where randomization is
used, preserve seeds, trial counts and raw sufficient statistics.
Resource limits are planning decisions, not evidence that a mathematical
solution cannot exist.

An output bundle should include:

\begin{itemize}
\tightlist
\item
  The statement being claimed, including all assumptions and
  quantifiers.
\item
  Construction or proof object, plus an independently runnable verifier
  and its trust assumptions.
\item
  Numerical/error ledger separating roundoff, discretization,
  truncation, optimization, sampling and modeling error where relevant.
\item
  Positive reference cases and negative controls that the verifier must
  reject.
\item
  Coverage: included domain, omitted domain, unresolved regions and
  termination reason.
\end{itemize}

These are requirements for future implementations, not a promise of
polynomial-size certificates for every problem.

\subsection{Checking specific certificate types}

\textbf{LP/SDP/SOS:} verify equality constraints, objective direction,
and cone feasibility using exact arithmetic or outward-rounded
intervals. Rationalizing a numerical solution can break PSD constraints.
Correct residuals with justified norm bounds; do not silently project
and reuse the old objective. Distinguish a proof of a polynomial
approximation from a proof for the original function.

\textbf{Interval numerics:} state function enclosures, domains and
outward rounding. A small root box needs existence and, if claimed,
uniqueness. Cover the complement for a global enumeration. Include
singularities and tails rather than discarding them numerically.

\textbf{SAT/SMT:} a solver's UNSAT label is not a certificate. Preserve
a proof supported by a checker and the proof that the encoding, symmetry
reduction and blocked solutions match the mathematical universe. DRAT
applies to an appropriate Boolean CNF proof, not automatically to every
real-arithmetic or SDP claim.

\textbf{Statistics:} specify the estimand, dependence assumptions,
confidence procedure and bias. Zero observed failures does not prove
zero failure probability. MCMC error bars based on estimated
autocorrelation are not generally rigorous mixing guarantees.

\subsection{Milestones and novelty}

Baseline means a working, independently validated result on a concrete
reference problem. Strong means a broader or sharper result with its
claimed evidence level. Research extension means a possible novel result
whose feasibility and significance need review. A negative theorem can
be valuable; an unsuccessful search can also be useful if its scope is
documented, but the two are not interchangeable.

Do not promise ``publication-quality'' in a fixed number of months.
Before claiming novelty, conduct a focused current literature comparison
against the exact theorem/model/metric. The sources in this revision
check key definitions and limitations; they are not an exhaustive
novelty audit across all 30 fields.

\end{document}
